\documentclass[11pt]{article}
\usepackage[margin=1in]{geometry}
\usepackage[T1]{fontenc}
\usepackage{lmodern,amsmath,amssymb,amsthm,microtype}
\usepackage[colorlinks=true,linkcolor=blue,citecolor=blue,urlcolor=blue]{hyperref}
\hypersetup{pdftitle={Representation constraints and polynomial-depth learning in Gaussian Bayesian networks},pdfauthor={Dmitry Vaintrob, Kaarel Hanni, OpenAI Codex}}
\newtheorem{theorem}{Theorem}[section]
\newtheorem{proposition}[theorem]{Proposition}
\newtheorem{lemma}[theorem]{Lemma}
\newtheorem{corollary}[theorem]{Corollary}
\newcommand{\op}{\mathrm{op}}
\newcommand{\poly}{\operatorname{poly}}
\title{Representation constraints and polynomial-depth learning\\
in Gaussian Bayesian networks}
\author{Dmitry Vaintrob \and Kaarel H\"anni \and OpenAI Codex}
\date{Research draft --- 19 October 2026}
\begin{document}
\maketitle
\begin{abstract}
We study Gaussian-prior tanh networks conditioned on a clean output
margin throughout the Boolean cube and on noiseless training signs.
For every fixed $\zeta>0$ and exponent $\gamma\ge1+\zeta$, a
posterior trained on $n$ examples at depth $L\ge3$ and sufficiently
large polynomial width assigns mass at least $1-\beta$ to Boolean
functions with same-depth unit-margin representatives whose product
of Frobenius norms is polynomial in $d,n$. A general polynomial
sample theorem would therefore require a corresponding approximation
representation theorem.

Conversely, a target with depth $D_0$ and product norm $P$ has
polynomial sample complexity, simultaneously at every width above
its actual bounded-weight network complexity, after only
$O(\log(d+D_0+P))$ additional learner layers. The original target
width and coefficient precision do not enter. We prove a
norm-preserving width compression theorem and a further global-margin
normal form whose entire Boolean description is polynomial in $P$,
independently of original depth and width, apart from specifying at
most $P^2$ relevant coordinate indices. At sufficiently large depth
and width, the minimum Gaussian energy above its depth baseline is
comparable to the logarithm of this intrinsic product complexity,
together with the input-normalization cost.

The positive results include every Boolean halfspace, independently
of its original geometric margin and coefficient precision, at every
learner depth $L\ge2$. We also prove polynomial sample bounds for
bounded-weight teachers with at most $K$ possible hidden states after
a fixed number of layers. The bound is polynomial in $K$, teacher
resources, and learner depth; no hidden-state separation is assumed.
For fan-in-two Boolean circuits of size $s$
and depth $h$, put $S=\max\{16,2+d+s\}$. At every learner depth
$L\ge\max\{2,h\}$, we obtain the upper bound
$S^{O_\zeta(1+h/\log\log S)}$ times a polynomial in learner depth and
statistical parameters. Apart from the cited Goldmann--Karpinski
threshold-to-majority simulation, the Hesse--Allender--Barrington
constant-depth threshold-circuit arithmetic, and Gaussian isoperimetry, we
give all analytic and representation arguments needed for these results. A sample bound jointly polynomial in width and depth for arbitrary
small deep circuits remains unresolved. A separate critical-prior
paper treats width-first learning at $\gamma=1$.
\end{abstract}

\section{Model and main conclusions}

The target $f:\{-1,1\}^d\to\{-1,1\}$ is learned from independent
examples $(x_i,f(x_i))$ drawn from an arbitrary input distribution $D$.
A learner of depth $L\ge2$ has $L-1$ hidden layers, each of width $N$,
and computes
\[
 h^0=(1,x),\qquad h^s=\tanh(W^sh^{s-1})\quad(1\le s<L),
 \qquad z_w(x)=v^Th^{L-1}.
\]
There are no biases after the augmented input. Its Gaussian energy is
\[
 \kappa_N(w)=\frac dN\|W^1\|_F^2+
    \sum_{s=2}^{L-1}\|W^s\|_F^2+N\|v\|_2^2.
\]
The raw Gaussian has independent entry variances
$t^2/d,t^2/N,t^2/N^2$, respectively. Put $\sigma^2=t^2/N$;
the exponent-$\gamma$ convention is $\sigma^2=N^{-\gamma}$.
In energy coordinates its density is proportional to
$\exp[-\kappa_N/(2\sigma^2)]$.

The posterior $Q_S$ is the raw Gaussian conditioned on
\[
 |z_w(x)|>1/2\quad\text{for every }x\in\{-1,1\}^d,
 \qquad f(x_i)z_w(x_i)>0\quad\text{for every training input}.
\]
``Unrestricted'' retains this all-input output-margin condition and
means that no additional internal robustness restriction is imposed.
Let $C_L(f)$ be the least width admitting a depth-$L$ representation
with all physical weights bounded by that width and signed output
greater than $1/2$ everywhere. This is the actual bounded-weight
network complexity at the specified learner depth.

Our new representation constraint is
Theorem~\ref{thm:posterior-energy-ceiling}. At every learner depth $L\ge3$, for every consistently labeled
sample of size at most $n$, every $0<\beta<1$, and all widths above a
polynomial threshold, posterior mass at least $1-\beta$ is supported on
Boolean functions admitting a unit-margin
depth-$L$ representative with product of Frobenius norms polynomial
in $d,n$. The threshold is polynomial in $d,n,L,\log(2/\beta)$, with degree depending
on the fixed gap $\zeta>0$ above exponent one. If $L$ is at least a
constant times $\log(d+n)$, one additional layer gives a representative
of energy $L+O(\log(d+n))$. Corollary~\ref{cor:approximation-energy-obstruction}
turns any function-level approximation-energy lower bound into an
actual posterior error lower bound. We do not establish such a barrier
for an explicit small deep target here.

Our new positive conclusion,
Theorem~\ref{thm:halfspace-unrestricted-all-depths}, covers every Boolean
halfspace with sample size polynomial in $d,L,\varepsilon^{-1}$ and
$\log(2/\delta)$, simultaneously for every $N\ge C_L(f)$ and every
$0<\sigma^2\le N^{-(1+\zeta)}$. Unlike a direct affine implementation,
the alternative shallow representation does not depend on the target's
geometric margin or original weight magnitudes. For halfspaces, the only
external representation ingredient is the classical Goldmann--Karpinski
threshold-to-majority simulation; the remaining analytic arguments are
included below.

Theorem~\ref{thm:fixed-prefix-finite-state} supplies a different
positive class: a bounded-weight teacher of depth $D$ whose state
after a fixed number of initial layers has at most $K$ values.
Its sample bound is polynomial in $K$, teacher resources, and learner
depth, at every $L\ge D$ and every actual width $N\ge C_L(f)$.
Only the teacher has a finite-state bottleneck; the prior receives no
additional restriction. For a bottleneck immediately after the first
layer, the bound at every learner depth $L\ge6$ is even independent of the
remaining computation.
Examples include first layers driven by a fixed number of groupwise
bit counts, or fixed-width first layers with polynomially bounded
integer weights. An arbitrary growing prefix depth or exponential
state count is not covered by this polynomial assertion.

The positive and negative directions meet at a function-level
representation question. Theorems~\ref{thm:product-target-log-padding}
and~\ref{thm:depth-independent-pf-normal-form} give, respectively,
a partial converse after logarithmic added depth and a finite Boolean
description controlled by product norm alone.
Theorem~\ref{thm:intrinsic-product-energy} then relates intrinsic
product complexity to minimum excess Gaussian energy at sufficiently
large depth and width. None of these statements identifies ordinary
circuit size with product complexity. For a general fan-in-two circuit
of size $s$ and depth $h$, put $S=\max\{16,2+d+s\}$. At every
$L\ge\max\{2,h\}$, the block construction gives
$S^{O_\zeta(1+h/\log\log S)}$ sample dependence, improving the direct
$S^{O_\zeta(h)}$ bound while retaining a gap from a joint polynomial.
Every fixed-depth threshold-circuit family is covered by a polynomial
bound; thus an explicit superpolynomial-sample counterexample within
$\mathsf{NC}^1$, at growing learner depth, would separate
$\mathsf{NC}^1$ from $\mathsf{TC}^0$.

Section~\ref{sec:pf-majority-closure} clarifies one possible route
from distribution-wise approximation to an exact representation.
Arbitrary added depth can be supplied at constant product cost.
Moreover, weak classifiers sharing one nonlinear feature map can be
combined at polynomial product cost. A general closure theorem for
classifiers with different feature computations remains open; simply
placing their networks in parallel can incur an exponential depth
factor. This separates an architecture-level obstruction from the
unresolved function-level representation question.

We measure error primarily by posterior Gibbs error
$\mathbb E_{w\sim Q_S}\Pr_D\{\operatorname{sign}z_w(X)\ne f(X)\}$.
If $p_S(x)$ is posterior probability of the correct label, then
\[
 D\{x:p_S(x)\le q\}\le
 \frac{\operatorname{GibbsErr}(Q_S)}{1-q}\qquad(0<q<1).
\]
Thus all positive results also cover correct-answer probability
thresholds $2/3$ and $1/2$, after constant changes in the Gibbs error
target. Posterior mixtures require no special assumption. The lower
bound in Corollary~\ref{cor:approximation-energy-obstruction} states
its separate consequence for the $2/3$ criterion explicitly.

The polynomial degrees may depend on $\zeta$; no uniform limit as
$\zeta\downarrow0$ is claimed. All results concern statistical sample
complexity, not efficient posterior computation. The draft has received
independent AI-agent proof audits, but is not Lean-certified.

\section{Changing the output margin at constant additional energy}
\label{sec:margin-transfer}

\begin{lemma}[One-layer margin transfer]
\label{lem:one-layer-margin-transfer}
Suppose a width-$N$, depth-$L$ tanh network with linear readout computes
a Boolean function $g$ with global signed margin at least $r>0$ and
Gaussian energy $K$. For every $r'>0$, there is a network at the same
width and depth $L+1$ computing $g$ with signed margin at least $r'$ and
energy at most
\[
 K+(r'/r)^2+(r')^2.
\]
The construction uses no biases after the original first layer.
\end{lemma}

\begin{proof}
Let $v$ be the old linear readout, and write
$c^2=N\|v\|_2^2$. The global margin implies $c\ge r>0$, because the
hidden coordinates have absolute value at most one. Define
\[
 a^2=\frac{r'}{cr},\qquad b=\frac{r'}{\tanh(ar)}.
\]
Replace the original readout by a hidden layer all of whose rows are
$av^T$, and give the new readout the constant entries $b/N$. The new
output is exactly $b\tanh(az(x))$. Oddness and monotonicity give the
claimed signed margin. The new two-layer energy is $a^2c^2+b^2$ in
place of the old readout energy $c^2$.

For every real $u$, $|\sinh u|\ge |u|$, so
\[
 \tanh^2u=\frac{\sinh^2u}{1+\sinh^2u}
       \ge\frac{u^2}{1+u^2}.
\]
Consequently
\[
 b^2\le(r')^2+\frac{(r')^2}{a^2r^2},\qquad
 a^2c^2+b^2\le 2c\frac{r'}r+(r')^2.
\]
Completing the square gives
$-c^2+2cr'/r\le(r'/r)^2$, proving the result.
\end{proof}

For example, a globally clean residual network with signed margin at
least $3/8$ can be converted to unit margin by adding one layer and at
most $73/9$ energy. The increase is independent of the old readout cost;
simply multiplying the old readout need not have that property.

Define the approximation energy
\[
 \mathcal E_{L,r}(f;\epsilon)
 =\inf\left\{\kappa_N(W):\begin{array}{l}
 N\ge1,\quad |z_W(x)|\ge r\text{ for all Boolean }x,\\
 \Pr_{X\sim D}\{\operatorname{sign}z_W(X)\ne f(X)\}<\epsilon
 \end{array}\right\},
\]
where the infimum ranges over depth-$L$ networks, $D$ is a specified
input distribution, and an empty infimum is
$+\infty$. The lemma implies the exact comparison
\[
 \mathcal E_{L+1,r'}(f;\epsilon)
 \le \mathcal E_{L,r}(f;\epsilon)+(r'/r)^2+(r')^2.
\]
Thus a unit-margin approximation-energy lower bound at one extra layer
implies a lower bound at any fixed smaller positive margin, losing only
a universal additive constant when the two margins are fixed.

\section{A posterior energy ceiling and a necessary representation condition}
\label{sec:posterior-energy-ceiling}

The following argument identifies a necessary condition for the desired
jointly polynomial unrestricted learning theorem. It concerns the actual
finite Gaussian scales above exponent one, not just a cold limit.
Throughout, the prior is conditioned on the original global clean margin
$|z(x)|>1/2$ and on the observed training signs.

\begin{lemma}[A uniform memorizing witness]
\label{lem:uniform-memorizing-witness}
For integers $d,n\ge1$, there are explicitly polynomial quantities
$M=M(d,n)$ and $T=T(d,n)\ge1$ such that every consistently labeled sample
of at most $n$ Boolean inputs admits a globally unit-margin interpolating
network at every width $N\ge M$ and every depth $L\ge3$, with energy at most
\[
 K=2+(L+14)(2T)^{1/(L-1)}.
\]
One may take
\[
 C=n+1,\quad
 R=\max\{4(n+1),(d+1)\log(16(n+1))\},
\]
\[
 M=\left\lceil\max\{4,C,2R,dC(d+1)R^2\}\right\rceil,
 \qquad T=MCR^2.
\]
In particular, if $L-1\ge\log(2T)$, then
\[
 K\le L+B(d,n),\qquad B(d,n)=16+32\log(2T)=O(\log(d+n)).
\]
\end{lemma}

\begin{proof}
Let $p_1,\ldots,p_m$ be the distinct negatively labeled sample points,
where $0\le m\le n$. Set $a=\log(16(n+1))$ and
\[
 q_j(x)=\frac{d-p_j\cdot x}{2},\qquad
 \tau_j(x)=\tanh\bigl(a(1/2-q_j(x))\bigr).
\]
Each $q_j$ is a nonnegative integer on the cube and is zero exactly at
$p_j$. Thus $\tau_j$ differs by at most $2e^{-a}$ from the sign which
equals $+1$ at $p_j$ and $-1$ elsewhere. Since the $p_j$ are distinct,
\[
 z(x)=-2\left[m-1+\sum_{j=1}^m\tau_j(x)\right]
\]
has the desired interpolating signs, equals $+2$ or $-2$ in the ideal
sign-feature limit, and differs from those values by at most
$4me^{-a}<1/4$. The formula also gives the constant $+2$ when $m=0$.
Implement its constant term using a first-layer constant feature $1/2$,
whose input bias is $\operatorname{atanh}(1/2)$. Pad with zero neurons
to $C=n+1$. Every physical weight is bounded by the stated $R$, and the
global signed margin is greater than $7/4$.

Zero-padding this shallow network to $M$ makes its first-layer energy
at most one and its weighted readout energy at most $T$. Balanced
cloning to any $N\ge M$ gives first energy at most two and readout
energy at most $2T$. The scalar-transmission construction of
Lemma~\ref{lem:scalar-transmission},
applied to this one-hidden-layer witness, gives the displayed bound
at every $L\ge3$.

For the final assertion, put $u=\log(2T)/(L-1)\le1$. The inequality
$e^u\le1+2u$ on $[0,1]$, and
$(L+14)/(L-1)\le16$, give
\[
 2+(L+14)e^u\le L+16+32\log(2T).
\]
\end{proof}

\begin{theorem}[A finite-scale posterior energy ceiling]
\label{thm:posterior-energy-ceiling}
Fix $d,n\ge1$, $L\ge3$, $\zeta>0$, and $0<\beta<1$. Use the quantities
$M,T,K$ above and put
\[
 Q_* = ((K+2)/L)^{L/2}.
\]
There is a width threshold polynomial in
$d,n,L,\log(2/\beta)$, with degree depending only on $\zeta$,
such that at every width above this threshold and every
$0<\sigma^2\le N^{-(1+\zeta)}$, the following holds for \emph{every}
consistently labeled sample of size at most $n$.

With posterior probability at least $1-\beta$, the posterior Boolean
function has a same-width, same-depth realization with global signed
margin greater than $3/8$ and energy at most $K+1$.
Rescaling its readout gives a unit-margin realization at the same
width and depth, with product of layer Frobenius norms bounded by
\[
 P_F^2\le d^{-1}e^{19}(2T)^{3/2}.
\]
Alternatively, it has a same-width realization at depth $L+1$, with
unit margin and energy at most $K+82/9$.

When $L-1\ge\log(2T)$, the last energy bound is
$L+O(\log(d+n))$.
\end{theorem}

\begin{proof}
Use Gaussian energy coordinates
$A_1=\sqrt{d/N}W^1$, $A_s=W^s$ internally, and $A_L=\sqrt N v$.
Their entries are independent $\mathcal N(0,\sigma^2)$ and their squared
Euclidean norm is $\kappa_N$. Write $t^2=N\sigma^2$.

We recall the explicit estimates from the Gaussian-distance bridge.
There is a closed, convex, centrally symmetric Gaussian set
$\mathcal T$ of probability at least $1/2$ whose normalized block
operator norms are at most $tb$, where
\[
 b^2=\max\{2(d+1)\log2+4\log8,
                    16\log9+8\log(8L)\}.
\]
Put
\[
 A=\left(\frac{K+2}{L-1}\right)^{(L-1)/2},\qquad
 \xi=\min\left\{1,\frac1{4\sqrt{L(K+1)}},
                         \frac1{8\sqrt2 LA}\right\}.
\]
If $tb\le\xi$, adding any $B\in\mathcal T$ to any network of energy
at most $K+1$ changes its output uniformly on the cube by at most $1/8$.
This follows by telescoping the layer perturbations and applying
arithmetic--geometric mean to the other $L-1$ layer norms.

For the realized sample, choose the unit-margin memorizing witness
$w_S$ above. Its shifted set $w_S+\mathcal T$ lies in the conditioning
event. Central symmetry and the Gaussian shift identity give
\[
 \mu(w_S+\mathcal T)\ge\tfrac12 e^{-K/(2\sigma^2)}.
\]
Gaussian isoperimetry gives
$\mu\{\operatorname{dist}_\kappa(w,\mathcal T)^2>K+1\}
\le e^{-(K+1)/(2\sigma^2)}$. Dividing yields
\[
 Q_S\{\operatorname{dist}_\kappa(w,\mathcal T)^2>K+1\}
       \le2e^{-1/(2\sigma^2)}.
\]
On the complementary event, subtract the metric projection of $w$
onto $\mathcal T$. The residual $U$ has energy at most $K+1$, and
its output differs by at most $1/8$ everywhere. The original clean
margin is greater than $1/2$, so the residual has the same Boolean
function and margin greater than $3/8$. Multiplying its readout by
$8/3$ gives unit margin without changing its depth.
Lemma~\ref{lem:one-layer-margin-transfer} adds at most $73/9$ energy
to obtain unit margin at one extra layer.

It suffices to take
\[
 N\ge\left\lceil\max\left\{M,2,(b/\xi)^{2/\zeta},
                                   2\log(2/\beta)\right\}\right\rceil.
\]
Indeed, $t\le N^{-\zeta/2}$ and $\sigma^{-2}\ge N$, so both estimates
hold simultaneously for every permitted scale. Since
$A\le\sqrt e Q_*$, this is a polynomial threshold in those quantities.
In fact, put $q^2=(2T)^{1/(L-1)}\ge1$. Since
$K+2\le(L+18)q^2$, one has for every $L\ge3$
\[
 Q_*^2\le e^{18}(2T)^{L/(L-1)}
          \le e^{18}(2T)^{3/2}.
\]
Thus the threshold is polynomial in $d,n,L,\log(2/\beta)$ for all
the stated depths. Since $K+1\le(L+17)q^2$, the same-depth residual
and arithmetic--geometric mean give, after rescaling the readout,
\[
 dP_F^2\le\frac{64}{9}\left(\frac{K+1}{L}\right)^L
 \le\frac{64}{9}e^{17}(2T)^{L/(L-1)}
 \le e^{19}(2T)^{3/2}.
\]
The additive-excess conclusion follows from the memorizing-witness
lemma. All constants and the set $\mathcal T$ are uniform over the
realized sample; only the translating witness depends on its labels.
\end{proof}

\begin{corollary}[Approximation-energy barriers force posterior error]
\label{cor:approximation-energy-obstruction}
Let $f$ and its input distribution $D$ be arbitrary. Under the conditions
of Theorem~\ref{thm:posterior-energy-ceiling}, suppose
\[
 \mathcal E_{L+1,1}(f;\epsilon)>K+82/9.
\]
Then, for every realized sample of at most $n$ examples labeled by $f$,
the posterior Gibbs error is at least $(1-\beta)\epsilon$.
The same conclusion holds if no globally unit-margin depth-$L$
network with $P_F^2\le d^{-1}e^{19}(2T)^{3/2}$ has $D$-error less
than $\epsilon$, allowing arbitrary finite widths.
If $p_S(x)$ is the posterior probability of the correct label, then
\[
 \Pr_D\{p_S(X)\le2/3\}
 \ge \max\left\{0,\frac{3(1-\beta)\epsilon-1}{2}\right\}.
\]
Conversely, if some sample gives posterior Gibbs error at most $r$,
then for every $r'>r/(1-\beta)$,
\[
 \mathcal E_{L+1,1}(f;r')\le K+82/9.
\]
There is also a globally unit-margin approximation of error less than
$r'$ at the original depth $L$, with
$P_F^2\le d^{-1}e^{19}(2T)^{3/2}$.
\end{corollary}

\begin{proof}
The energy ceiling puts posterior mass at least $1-\beta$ on Boolean
functions with unit-margin representatives of energy at most $K+82/9$.
Under the barrier, every such function has error at least $\epsilon$.
The same posterior event has the same-depth product-norm realization
in the theorem, proving the alternative barrier statement.
For the confidence claim, if the displayed confidence-error probability
is $e$, then Gibbs error is at most $(1-e)/3+e$.
For the converse, condition the posterior on its event of mass at least
$1-\beta$. Its average error is at most $r/(1-\beta)$, so for every
larger $r'$ it contains a function of error less than $r'$.
\end{proof}

Thus a jointly polynomial sample theorem for all small deep targets
would imply a representation theorem at every $L\ge3$: every target
learned with polynomially many samples must admit a constant-error,
unit-margin approximation at the same depth $L$ whose product of layer
Frobenius norms is polynomial in the target parameters and $d$.
At depths sufficiently large
compared with $\log(d+n)$, every such target must admit a constant-error
approximation of energy $L+O(\log(d+n))$ at depth $L+1$. In particular,
if $n$ is polynomial in $d,C,L$, the required excess energy is only
$O(\log(d+C+L))$. This conclusion does not depend on choosing the
product-norm proof strategy for the upper bound; it is forced by the
posterior itself. No example of a small target violating this necessary
condition has been established here.

\section{Boolean halfspaces need no geometric-margin assumption}
\label{sec:halfspaces-no-geometric-margin}

The shallow-target theorem applies to every Boolean linear threshold
function, even when its most natural affine representation has extremely
small geometric margin. The relevant fact is a classical circuit
simulation theorem; retaining the original affine representation would
miss this conclusion.

\begin{lemma}[Classical threshold-to-majority simulation]
\label{lem:goldmann-karpinski-simulation}
There is a universal polynomial $p$ with the following property.
Every Boolean threshold function $f:\{-1,1\}^d\to\{-1,1\}$ can be
written
\[
 f(x)=\operatorname{sign}\!\left(b_0+
       \sum_{j=1}^m b_j\operatorname{sign}
       \left(a_{j0}+\sum_{i=1}^d a_{ji}x_i\right)\right),
\]
where $m\le p(d)$, all $a_{ji},b_j$ are integers of absolute value
at most $p(d)$, and every displayed sign has a nonzero argument on
every Boolean input.
\end{lemma}

This is the depth-two majority simulation of
Goldmann and Karpinski~\cite{goldmann1998}. Majority gates allow repeated and negated
inputs and constants; collecting these wires gives the integer
coefficients in the lemma. Their size convention counts wires, so their
polynomial size bound also bounds all these coefficients. Their paper
also gives the statement directly in terms of polynomially bounded
integer weights. This is an external circuit theorem, not a new proof
of threshold simulation here. A threshold function initially specified
with a fixed convention for ties can be put in strictly separating
form by an arbitrarily small bias change on the finite Boolean cube.

\begin{lemma}[A polynomial shallow tanh realization]
\label{lem:halfspace-shallow-tanh}
Every Boolean threshold function on $d$ inputs has a one-hidden-layer
tanh network with linear readout, global signed margin at least $3/2$,
and both width and maximum absolute physical weight bounded by a
universal polynomial in $d$. These bounds are independent of the
coefficients and geometric margin of its original threshold
representation.
\end{lemma}

\begin{proof}
Use Lemma~\ref{lem:goldmann-karpinski-simulation}. Put
$A=\max\{1,\sum_{j=1}^m|b_j|\}$ and
$\lambda=\tfrac12\log(8A)$. Each inner affine form is a nonzero
integer, so, with $h_j(x)$ denoting its sign,
\[
 \left|\tanh\!\left(\lambda
      \left(a_{j0}+\sum_i a_{ji}x_i\right)\right)-h_j(x)\right|
       \le 2e^{-2\lambda}=\frac1{4A}.
\]
Consequently replacing all bottom signs by these tanh features changes
the top affine form by at most $1/4$. Its exact signed value is a
positive integer and hence at least one. Therefore
\[
 z(x)=2b_0+2\sum_{j=1}^m b_j
       \tanh\!\left(\lambda
       \left(a_{j0}+\sum_i a_{ji}x_i\right)\right)
\]
satisfies $f(x)z(x)\ge3/2$ everywhere.

Implement the constant $2b_0$ by one additional hidden feature equal
to $1/2$, using first-layer bias $\operatorname{atanh}(1/2)$ and
readout coefficient $4b_0$. The resulting width is $m+1$, and every
physical weight has magnitude at most
\[
 \max\{\lambda p(d),4p(d),\operatorname{atanh}(1/2)\},
\]
which is polynomial in $d$. No later bias is used.
\end{proof}

\begin{theorem}[Unrestricted learning of every Boolean halfspace]
\label{thm:halfspace-unrestricted-all-depths}
Fix $\zeta>0$. There is a polynomial $P_\zeta$ such that, for every
Boolean threshold target $f$ on $d$ inputs, every learner depth
$L\ge2$, and every input distribution, a sample of size
\[
 n\ge P_\zeta\bigl(d,L,\varepsilon^{-1},\log(2/\delta)\bigr)
\]
has, with probability at least $1-\delta$, the following simultaneous
guarantee. At every width $N\ge C_L(f)$ and every prior scale
$0<\sigma^2\le N^{-(1+\zeta)}$, the original unrestricted Gaussian
posterior has Gibbs error at most $\varepsilon$. Here $C_L(f)$ is
the actual bounded-weight network complexity at the specified learner
depth, and the posterior is conditioned only on the original global
clean output margin and the observed training labels.

In particular, the guarantee holds for every $\gamma\ge1+\zeta$,
with no dependence on an inverse geometric margin or on the precision
of a specified affine representation of $f$.
\end{theorem}

\begin{proof}
Lemma~\ref{lem:halfspace-shallow-tanh} supplies polynomial-in-$d$
values of the shallow width $C$ and physical weight bound $R$ in
Theorem~\ref{thm:shallow-all-depths}. Substitute these values in that
theorem. Its smaller-width argument already covers every
$C_L(f)\le N<M$ below the polynomial-width constructed witness; no
identification of $C_L(f)$ with that constructed width is required.
\end{proof}

This is an information-theoretic posterior guarantee. It does not
assert that sampling the conditioned posterior, or recovering the
alternative shallow representation from examples, is computationally
efficient.

\section{Targets with small Fourier spectral norm}
\label{sec:fourier-spectral-targets}

For a Boolean function $f:\{-1,1\}^d\to\{-1,1\}$, write
\[
 f(x)=\sum_{S\subseteq[d]}\widehat f(S)\chi_S(x),
 \qquad \chi_S(x)=\prod_{i\in S}x_i,
 \qquad A(f)=\sum_S|\widehat f(S)|.
\]
The Fourier coefficients here use the uniform distribution, even when
the learning distribution is different. Parseval implies $A(f)\ge1$.
The following sufficient condition imposes neither a Fourier-degree
bound nor a bound on the number of nonzero Fourier coefficients.

\begin{lemma}[A shallow realization from Fourier spectral norm]
\label{lem:fourier-spectral-shallow}
Every such $f$, with $A=A(f)$, has a one-hidden-layer tanh realization
with linear readout, signed margin at least $13/8$, width at most
\[
 1+d\left\lceil128A^2(d+2)\log2\right\rceil,
\]
and physical weights bounded by
\[
 \max\{3A,\tfrac d2\log(32Ad),1\}.
\]
\end{lemma}

\begin{proof}
Sample a set $S$ with probability $|\widehat f(S)|/A$ and attach
the sign $s_S=\operatorname{sign}\widehat f(S)$.
The random function $A s_S\chi_S$ has mean $f$ and takes values
in $[-A,A]$. For
$m=\lceil128A^2(d+2)\log2\rceil$, Hoeffding's inequality and a
union bound over the $2^d$ inputs show that, with positive probability,
\[
 \left|\frac A m\sum_{j=1}^m s_{S_j}\chi_{S_j}(x)-f(x)\right|
 \le\frac18\qquad\text{for every }x.
\]
Fix one such choice.

Each parity on $k\ge1$ variables admits an elementary real-output
threshold expansion. Put $t=\sum_{i\in S}x_i$ and
$\theta_j=-k+2j-1$ for $1\le j\le k$. On the permitted values
$t\in\{-k,-k+2,\ldots,k\}$,
\[
 \chi_S(x)=b_k+\sum_{j=1}^k(-1)^{k-j}
                  \operatorname{sign}(t-\theta_j),
 \qquad b_k\in\{0,1\}.
\]
Indeed, both sides start at $(-1)^k$ and have the same jump at
each threshold; explicitly $b_k=1$ for even $k$ and zero for odd $k$.
Every threshold has absolute affine margin at least one.
With $\lambda=\tfrac12\log(32Ad)$, replace each sign by
$\tanh(\lambda(t-\theta_j))$. The resulting parity approximation
has uniform error at most
\[
 2k e^{-2\lambda}\le\frac1{16A}.
\]
The empty parity is the constant one. A single constant hidden neuron,
whose output is $\tanh1$, supplies all constant terms.

Replacing the $m$ sampled parities by these approximations adds error
at most $1/16$. Thus the resulting scalar output is within $3/16$ of
$f$ everywhere. Doubling the readout gives signed margin at least
$13/8$. There are at most $md$ nonconstant hidden neurons. Their
first-layer coefficients and biases have magnitude at most $d\lambda$;
their output coefficients have magnitude at most $2A/m$. The single
constant neuron's output coefficient has magnitude at most
$2A/\tanh1<3A$. These are the stated bounds.
\end{proof}

\begin{corollary}[Unrestricted learning from small spectral norm]
\label{cor:fourier-spectral-unrestricted}
For each fixed $\zeta>0$, there is a polynomial $P_\zeta$ such that
for every Boolean target $f$, every input distribution, and every
learner depth $L\ge2$, a sample size
\[
 n\ge P_\zeta\bigl(d,A(f),L,\varepsilon^{-1},\log(2/\delta)\bigr)
\]
gives, with probability at least $1-\delta$, Gibbs error at most
$\varepsilon$ for the original unrestricted Gaussian posterior,
simultaneously for every $N\ge C_L(f)$ and every
$0<\sigma^2\le N^{-(1+\zeta)}$.
The posterior retains its original conditioning on a global clean
output margin and the training labels.
\end{corollary}

\begin{proof}
Substitute Lemma~\ref{lem:fourier-spectral-shallow} into
Theorem~\ref{thm:shallow-all-depths}. All witness parameters are
polynomial in $d$ and $A(f)$. The all-width conclusion is the
all-width conclusion of that theorem and does not require
$C_L(f)$ to equal the width of the constructed shallow realization.
\end{proof}

Both the realization lemma and its learning corollary also hold with
$A(f)$ replaced by $\sum_S|\widehat g(S)|$ for any real-valued Fourier
polynomial $g$ satisfying $f(x)g(x)\ge1$ on every input. The same
sampling argument approximates $g$ to error $3/16$, which still leaves
the stated signed margin after doubling the readout. Thus a small
normalized coefficient sum for a margin-separated vote of parities
is enough, even when the Boolean sign function itself has a larger
Fourier norm.

For example, every parity has $A(f)=1$, independently of its degree.
More generally, a function with $s$ nonzero Fourier coefficients has
$A(f)\le\sqrt s$ by Cauchy--Schwarz and Parseval. These are sufficient
representation conditions. The argument does not assert that all small
Boolean formulas or all narrow branching programs have polynomial
Fourier spectral norm.

\begin{corollary}[Two-state branching programs]
\label{cor:two-state-branching}
Every Boolean function computed by a deterministic layered width-two
branching program of length $T$ has $A(f)\le T+1$. Consequently the sample bound
in Corollary~\ref{cor:fourier-spectral-unrestricted} is polynomial in
$d,T,L,\varepsilon^{-1},\log(2/\delta)$ for this class.
The program may query variables repeatedly, and its two nodes in a
layer may query different variables.
\end{corollary}

\begin{proof}
Encode the two states as $-1,+1$, and let $s_j(x)$ be the state after
$j$ transitions. Given the old state $+1$ or $-1$, the next state is
a Boolean function $u_+(x)$ or $u_-(x)$ of at most one input bit.
Each such function is constant or a signed input bit, so
$A(u_+)=A(u_-)=1$. The transition can be written
\[
 s_{j+1}(x)=a_j(x)s_j(x)+b_j(x),\qquad
 a_j=(u_+-u_-)/2,\quad b_j=(u_++u_-)/2.
\]
The Fourier coefficient norm is subadditive and submultiplicative:
the latter follows because multiplication of parities uses symmetric
difference of their index sets. Thus $A(a_j),A(b_j)\le1$ and
\[
 A(s_{j+1})\le A(s_j)+1.
\]
The initial state is constant, so $A(s_0)=1$. The final output is
either constant or a signed copy of the final state, proving the claim.
Unused state slots can be filled arbitrarily. If the program uses
absorbing states for early termination, these states count toward its
width-two bound.
\end{proof}

The same proof allows a node to query a parity rather than a single
bit, since a signed parity also has Fourier norm one. Decision lists
with $T$ parity tests likewise have Fourier norm at most $T+1$, by a
separate direct argument: write the output as the sum of its $T+1$
signed first-accept and default indicators. Each indicator is a
product of factors $(1\pm\chi_S)/2$, each of Fourier norm at most
one, so submultiplicativity and the triangle inequality give the
bound. These extensions concern representation and sample bounds
only; no efficient algorithm for identifying a program from samples
is asserted.

\section{Constant-depth threshold circuits and a lower-bound barrier}
\label{sec:threshold-circuits}

The coefficient-independent halfspace construction also clarifies which
explicit targets could disprove the general unrestricted sample
conjecture. Here a threshold circuit may have arbitrary real weights,
and its size counts wires, including input and constant wires.

\begin{lemma}[Threshold circuits as bounded-weight tanh networks]
\label{lem:threshold-circuit-tanh}
There is a universal constant $a$ such that every depth-$h$, $h\ge1$, threshold
circuit of wire size at most $s$ on $d$ Boolean inputs has a tanh
realization with linear readout, unit global margin, depth at most
$2h$, and both width and maximum absolute physical weight bounded by
$(2+d+s)^a$. It uses no biases after the augmented input. The bound is
independent of the original real coefficients and of every original
gate's geometric margin.
\end{lemma}

\begin{proof}
Replace each threshold gate by its depth-two polynomial-size,
polynomial-integer-weight simulation from
Lemma~\ref{lem:goldmann-karpinski-simulation}. Each gate has at most
$s$ inputs. The total wire count, coefficient bound, and maximum
weighted fan-in of the resulting circuit are bounded by a universal
polynomial in $d+s$. Its depth is at most $2h$. Subdivide skip wires by
identity threshold gates to make the circuit layered; since $h\le s$,
this costs only another polynomial factor. Pad with identity gates as
needed so that the output is at a common depth $q=2h\ge2$.

Every gate can be taken to have integer affine form with absolute value
at least one on Boolean inputs. Let $B\ge1$ bound the sum of the
absolute coefficients of any such affine form, including its constant.
Put $\delta=1/(4B)$ and $\lambda=2\log(8B)$. Replace every gate except
the final output gate by tanh applied to $\lambda$ times its affine
form. If all preceding features approximate their intended signs to
accuracy $\delta$, the affine perturbation is at most $B\delta=1/4$.
Its signed value is therefore at least $3/4$, and the new feature error
is at most
\[
 2e^{-3\lambda/2}=\frac{2}{(8B)^3}\le\delta.
\]
This inductive accuracy bound does not deteriorate with depth.

To implement constants without later biases, add a constant feature
equal to $1/2$ at every hidden layer. The first such feature has bias
$\operatorname{atanh}(1/2)$; each later one receives weight
$2\operatorname{atanh}(1/2)$ from the preceding constant feature.
An affine constant $b$ is then supplied by weight $2\lambda b$ to
this feature. Thus all constants are exact. The last threshold is
replaced by twice its affine form, as a linear readout; its signed
value is at least $3/2$. Uniform zero-padding of the hidden layers
preserves the function. Every used parameter is polynomial in $d+s$,
as asserted.
\end{proof}

\begin{corollary}[Precision-independent threshold-circuit learning]
\label{cor:threshold-circuit-learning}
Fix $\zeta>0$. A target with a depth-$h$, wire-size-$s$ threshold circuit
has unrestricted Gaussian posterior sample complexity
\[
 (2+d+s)^{O_\zeta(h)}
 \bigl(L+\varepsilon^{-1}+\log(2/\delta)\bigr)^{O_\zeta(1)}
\]
at every learner depth $L\ge2h$, simultaneously for every
$N\ge C_L(f)$ and every $0<\sigma^2\le N^{-(1+\zeta)}$.
No original coefficient magnitude or geometric margin enters the bound.
In particular, for each fixed $h$, this is polynomial in the circuit
size, input dimension, learner depth, and statistical parameters.
\end{corollary}

\begin{proof}
Apply the fixed-intrinsic-depth theorem and its quantitative corollary
below to Lemma~\ref{lem:threshold-circuit-tanh}, with intrinsic depth
$D_0=2h$ and polynomial bounds on the tanh width and weights.
\end{proof}

Recall that nonuniform $\mathsf{TC}^0$ consists of Boolean families
having polynomial-size threshold circuits of some fixed depth.
Consequently every such family enjoys a polynomial sample guarantee
under the present posterior at all sufficiently large learner depths.
This gives the following precise limitation on a negative-result route.

\begin{corollary}[A circuit-complexity consequence of a counterexample]
\label{cor:tc0-counterexample-barrier}
Fix $\zeta>0$. Suppose a family $(f_d)$, at learner depths $L(d)\to\infty$,
requires more than every polynomial in $d,L(d)$ examples for some fixed
positive error and confidence targets, under the present unrestricted
posterior and the all-admissible-width guarantee with
$\gamma\ge1+\zeta$. Then $(f_d)\notin\mathsf{TC}^0$.
In particular, a family in $\mathsf{NC}^1$ with this property would
separate $\mathsf{NC}^1$ from $\mathsf{TC}^0$.
\end{corollary}

\begin{proof}
If the family were in $\mathsf{TC}^0$, it would have fixed depth $h$
and wire size polynomial in $d$. Eventually $L(d)\ge2h$, so
Corollary~\ref{cor:threshold-circuit-learning} would supply a polynomial
sample bound simultaneously for every admissible width and stated
Gaussian scale, a contradiction.
\end{proof}

This is a logical consequence, not a claimed circuit separation or a
proof that unrestricted learning is always polynomial. It also does not
cover exponents approaching one with dimension, fixed learner depths
below the alternative representation's depth, or the mean-field case.

\section{Finite hidden-state bottlenecks give polynomial depth bounds}
\label{sec:finite-state-bottleneck}

Boolean inputs permit a positive result that has no analogue for a
continuously varying scalar input. If a short initial part of the teacher
produces only polynomially many distinct hidden states, its remaining
computation can be replaced by a lookup circuit. The remaining teacher
layers may be arbitrarily expansive. No additional condition is imposed
on the learner's prior.

\begin{proposition}[A finite first-layer state set]
\label{prop:first-layer-finite-state}
Suppose a tanh network in the present architecture computes
$f:\{-1,1\}^d\to\{-1,1\}$ and its first hidden vector has $C$
coordinates and at most $K$ distinct values on the cube. Then $f$ has a
depth-three threshold circuit with $O(KC(d+1))$ wires. Consequently it
has a globally unit-margin tanh realization of depth six, with width
and individual weights polynomial in $d,C,K$.
There is no bound on the depth, coefficients, or numerical stability
of the specified teacher's remaining computation.
\end{proposition}

\begin{proof}
For every attained hidden vector $v$, its label is well-defined because
all subsequent computation depends only on that vector. Tanh is
injective. Thus the condition that its $j$th coordinate equals $v_j$
can be tested by requiring the original affine preactivation to lie
between two thresholds that isolate its corresponding value from all
other attained preactivation values. There are only finitely many such
values. Choose separating endpoints strictly between consecutive values;
at an extreme, omit the unnecessary test or use a constant predicate.

Equality to $v$ is therefore a conjunction of at most $2C$ affine
threshold predicates of the original input. Use one threshold layer for
these predicates, a second for each state conjunction, and a final OR
over the positively labeled states. There are at most $K$ states and
each first-layer predicate uses at most $d+1$ wires. Constant functions
are immediate. Lemma~\ref{lem:threshold-circuit-tanh} gives the stated
bounded-weight tanh realization. In particular, arbitrarily small gaps
between attained real values cause no coefficient-size loss in the
resulting realization.
\end{proof}

The proposition already gives a sample polynomial in $d,C,K,L$ and the
statistical parameters at every learner depth $L\ge6$, simultaneously
for all $N\ge C_L(f)$ and all
$0<\sigma^2\le N^{-(1+\zeta)}$, for each fixed $\zeta>0$.
The next result allows a bottleneck after any fixed number of teacher
layers. Its proof uses binary precision, rather than a lower bound on
the gaps between all attained states.

\begin{lemma}[Fixed-depth tanh evaluation by threshold circuits]
\label{lem:fixed-prefix-tc-evaluation}
Fix $j\ge1$. Consider the first $j$ hidden layers of a tanh network
with widths at most $C$, input $(1,x)$ for $x\in\{-1,1\}^d$, and
weights of magnitude at most $B\ge1$. For every integer $p\ge1$, a
threshold circuit of depth depending only on $j$ and wire size
polynomial in $d,C,B,p$ produces a dyadic vector $\widehat h(x)$ such
that
\[
 \sup_x\|\widehat h(x)-h^j(x)\|_\infty\le2^{-p}.
\]
The polynomial degree may depend on $j$. Its advice consists of
polynomially many bits of the specified real parameters; no
computability assumption on those real parameters is required.
\end{lemma}

\begin{proof}
We use the theorem of Hesse--Allender--Barrington
\cite{hesse2002}: division and iterated multiplication of polynomially
many polynomial-bit integers have polynomial-size constant-depth
threshold circuits. Addition, iterated addition, and comparison also
have such circuits. Only their nonuniform consequence is needed here.

First consider a dyadic scalar $t=a/2^b$ with $|t|\le R$, $R\ge1$.
Tanh can be approximated to error $2^{-p}$ by a constant-depth
threshold circuit of size polynomial in $b,R,p$. Indeed, take the
degree-$m$ Taylor polynomial of $e^{2t}$, with
$m=\lceil40(R+p+2)\rceil$. Its uniform remainder is at most
\[
 e^{2R}\frac{(2R)^{m+1}}{(m+1)!}
 \le e^{2R}\left(\frac{2eR}{m+1}\right)^{m+1}
 <2^{-p-4}.
\]
The polynomial is the exact rational number $U/V$, where
\[
 V=2^{bm}m!,\qquad
 U=\sum_{k=0}^m (2a)^k2^{b(m-k)}\frac{m!}{k!}.
\]
All products and sums involve polynomially many integers with
$O(m(b+\log(R+1)+\log(m+1)))$ bits. Negative summands are handled
exactly; cancellation does not introduce numerical error.
The quoted arithmetic theorem computes these quantities at constant
depth. Since $U/V$ is within $2^{-p-4}$ of a nonnegative number,
$U/V+1\ge3/4$. Hence
\[
 \frac{U-V}{U+V}
\]
approximates $(e^{2t}-1)/(e^{2t}+1)=\tanh t$ to error less than
$2^{-p-2}$. A final rational division and dyadic rounding give the
required error. Clamping the result to $[-1,1]$ can only decrease its
distance from $\tanh t$.

To evaluate the prefix, round each parameter to absolute error
$2^{-q}$ and approximate every tanh to the same error, where
\[
 q=p+\left\lceil\log_2\bigl(4j(d+C+2)(2CB)^j\bigr)\right\rceil.
\]
Clamp all approximate hidden coordinates to $[-1,1]$.
If $E_s$ is the maximum hidden-coordinate error after layer $s$, then
\[
 E_1\le(d+2)2^{-q},\qquad
 E_s\le CB E_{s-1}+(C+1)2^{-q}\quad(2\le s\le j).
\]
Therefore
\[
 E_j\le j(d+C+2)(2CB)^j2^{-q}\le2^{-p}.
\]
The rounded preactivations have magnitude at most
$(d+C+1)(B+1)$ and polynomial bit length. Every affine calculation
uses only a constant number of the arithmetic operations above, and
each tanh uses the preceding scalar construction. Composing the fixed
number $j$ of stages preserves constant threshold depth and polynomial
wire size. The rounded matrices have polynomial total description
length. Their bits are fixed circuit advice, rather than quantities
the circuit is required to extract from arbitrary real numbers.
\end{proof}

\begin{theorem}[A finite bottleneck after a fixed prefix]
\label{thm:fixed-prefix-finite-state}
Fix $j\ge1$ and $\zeta>0$. Suppose a width-$C$, depth-$D>j$ tanh
teacher has weights bounded by $B\ge1$, signed output at least $1/2$
on the Boolean cube, and at most $K$ distinct vectors after its $j$th
hidden layer. There are a constant $q_j$ and a polynomial $p_j$, with
degrees depending only on $j$, such that $f$ has a globally
unit-margin tanh realization of depth $q_j$, width and weights at most
$p_j(d,C,B,K,D)$.

Consequently the original unrestricted Gaussian posterior has sample
complexity polynomial in
\[
 d,C,B,K,D,L,\varepsilon^{-1},\log(2/\delta)
\]
at every learner depth $L\ge q_j$, simultaneously for every
$N\ge C_L(f)$ and every $0<\sigma^2\le N^{-(1+\zeta)}$.
The same polynomial form holds at every learner depth $L\ge D$,
including the possible bounded-depth range below $q_j$.
The constants and polynomial degrees may depend on $j$ and $\zeta$,
but not on the remaining teacher depth $D-j$.
\end{theorem}

\begin{proof}
In the infinity norm, the teacher suffix from its $j$th hidden layer
to the linear readout is Lipschitz with constant at most
\[
 S=(CB)^{D-j}.
\]
Oppositely labeled inputs have outputs differing by at least one.
Consequently their $j$th hidden states are separated in infinity norm
by at least $S^{-1}$. Set
\[
 p=\lceil\log_2(32S)\rceil,\qquad \eta=2^{-p}.
\]
The required precision is only $O(1+D\log(2CB))$ bits, despite the
possibly exponentially small separation.

Apply Lemma~\ref{lem:fixed-prefix-tc-evaluation} to produce
$\widehat h(x)$ within $\eta$ of the true prefix state. For each
attained state $v$, hardcode a dyadic center $\widehat v$ within
$\eta$ of $v$, and hardcode its Boolean label. Output $+1$ exactly
when $\widehat h(x)$ lies within $3\eta$ in every coordinate of at
least one positive center. A positive input lies within $2\eta$ of
its own center. A negative input is at distance at least
$S^{-1}-2\eta\ge30\eta$ from every positive center. Thus the
lookup returns the correct label everywhere.

There are at most $K$ centers, each described by $O(Cp)$ bits.
Subtraction, comparison, a conjunction over coordinates, and a
disjunction over centers are constant-depth threshold operations of
polynomial size. Together with the fixed-depth prefix evaluation,
this gives a threshold circuit of depth depending only on $j$ and
size polynomial in $d,C,B,K,D$. Apply
Lemma~\ref{lem:threshold-circuit-tanh} and pad the resulting fixed
depth to a constant $q_j$. This proves the representation statement.

The fixed-intrinsic-depth learning theorem gives the claimed
all-admissible-width guarantee for $L\ge q_j$. To cover every
$L\ge D$, it remains only to consider $D<q_j$ and $L<q_j$.
There are finitely many such intrinsic teacher depths, bounded in terms
of $j$ alone. Double the supplied teacher's readout and apply the same
fixed-intrinsic-depth theorem directly to that teacher. Taking the
maximum of these finitely many polynomial bounds gives the stated
form. The prior throughout retains only its original global output
margin restriction and the observed training signs.
\end{proof}

In the representation proof it would suffice to cover the attained
prefix states by $K$ infinity-norm balls of radius $\eta$, with centers
chosen among those states. Use boxes of radius $4\eta$ in the lookup.
Opposite labels cannot occur in the same such ball, and the same
separation argument applies. Exact finite-state count is a particularly
simple sufficient condition; no separation of same-label states is
needed.

\begin{corollary}[Discrete input statistics]
\label{cor:discrete-statistic-bottleneck}
The following teachers satisfy Proposition~\ref{prop:first-layer-finite-state}.
\begin{enumerate}
\item If the nonconstant entries of first-layer row $j$ are
$A_{ji}/Q_j$, with integers $A_{ji}$ and $Q_j>0$, put
$R_j=\sum_{i=1}^d|A_{ji}|$. Its state count is at most
\[
 K\le\prod_{j=1}^C(2R_j+1).
\]
Its first-layer biases may be arbitrary real numbers. In particular,
first-layer input weights bounded by $B$ with common denominator $Q$
give $K\le(2dBQ+1)^C$. For fixed $C$ and polynomially bounded $B,Q$,
the sample bound is polynomial in input dimension and learner depth,
independently of all later teacher computation.
\item If the first-layer matrix has at most $k$ distinct nonconstant
input-column vectors, let $d_1,\ldots,d_k$ be their multiplicities.
Then
\[
 K\le\prod_{a=1}^k(d_a+1)\le(d+1)^k.
\]
For fixed $k$, the sample bound is polynomial in $d,C,L$ and the
statistical parameters, with no dependence on the precision or
magnitude of the specified teacher coefficients. Functions of a fixed
number of coordinate-block Hamming weights are examples.
\end{enumerate}
\end{corollary}

\begin{proof}
In the first case, the integer sum in row $j$ lies between $-R_j$
and $R_j$, so it has at most $2R_j+1$ values. A fixed bias and tanh
do not increase this count. In the second case, the first affine map
depends only on the sum of the Boolean inputs in each column group.
A group of size $d_a$ has exactly $d_a+1$ possible sums. All later
computation factors through the resulting state vector. Apply the
proposition and the fixed-intrinsic-depth learning theorem.
\end{proof}

These results use finite information passed by a teacher prefix,
rather than stability of the prior's networks. They do not imply a
joint polynomial theorem for an arbitrary growing prefix depth or
unrestricted state count. Even one bounded real affine projection can
take $2^d$ distinct values on the cube. Rounding that projection while
preserving a long suffix's output signs may require exponentially fine
precision in depth, so it does not establish the required polynomial
state bound for a general teacher.

\section{A depth-blocking bound for general Boolean circuits}
\label{sec:boolean-blocking}

A simple circuit transformation slightly improves the exponential
dependence on growing intrinsic depth. It does not make the dependence
polynomial for general small circuits.

\begin{proposition}[Blocking Boolean layers]
\label{prop:boolean-layer-blocking}
Let $f$ be computed by an AND/OR/NOT circuit of fan-in at most two,
depth $h\ge1$, and at most $s$ gates and input wires. Put
\begin{align*}
 S&=\max\{16,2+d+s\},\qquad
 b=\max\{1,\lfloor\log_2\log_2 S\rfloor\},\\
 a&=\lfloor h/b\rfloor,\qquad r=h-ab,\qquad
 q=\max\{2,2a+\min\{2,r\}\}.
\end{align*}
Then $f$ has a depth-$q$ tanh representation with unit global margin,
width and physical weights bounded by a universal polynomial in $S$.
Consequently, for each fixed $\zeta>0$, at every learner depth $L\ge q$
the unrestricted posterior has a sample bound of the form
\[
 S^{\,O_\zeta(1+h/\log\log S)}
 \bigl(L+\varepsilon^{-1}+\log(2/\delta)\bigr)^{O_\zeta(1)},
\]
uniformly for every actual width $N\ge C_L(f)$ and every
$0<\sigma^2\le N^{-(1+\zeta)}$.
One has $q\le\max\{2,h\}$, so in particular the result applies
at every learner depth $L\ge\max\{2,h\}$.
For polynomial-size logarithmic-depth circuits, this is
$\exp[O_\zeta(\log^2 S/\log\log S)]$ times a polynomial in the
learner depth and statistical parameters.
\end{proposition}

\begin{proof}
First insert identity gates along skip wires and at an output of depth
less than $h$. The resulting layered circuit has polynomial size in $S$,
since $h\le s$ after removing unused gates. Divide its layers into
consecutive blocks of at most $b$ layers.

Each gate at the end of a block depends on at most $2^b$ signals at the
beginning of that block. Regard its subcomputation as a Boolean function
of these boundary signals, and express that function by its full DNF.
The DNF has at most
\[
 2^{2^b}\le S
\]
terms, each of length at most $2^b\le\log_2 S$. Implementing its
conjunctions in one threshold layer and its final disjunction in the
next uses polynomially many wires and integer coefficients bounded by
a polynomial in $S$. A gate whose function is constant can be
implemented by a constant threshold. Use two threshold layers for
every full block. A final one-layer block can simply retain its
original AND/OR/NOT threshold gates, rather than be expanded to two
layers; a longer final partial block uses two layers. The resulting
threshold depth is $2a+\min\{2,r\}\le h$. When $h=1$, pad to depth
two by inserting one identity threshold layer. All blocks together
therefore form a depth-$q$ threshold circuit of polynomial size and
polynomial integer weights. Choose odd integer thresholds after doubling literal
coefficients, so every gate has signed preactivation at least one on
Boolean inputs.

Here the threshold circuit already has small integer weights; no
threshold-to-majority replacement is needed. Convert its first $q-1$
layers directly to tanh. If $B\ge1$ bounds each gate's total absolute
incoming coefficients, including its bias, use gain
$\lambda=2\log(8B)$ and accuracy $1/(4B)$ at every hidden layer.
Inductively the affine perturbation from preceding approximate signs
is at most $1/4$, and the tanh error is at most
$2e^{-3\lambda/2}\le1/(4B)$. Biases are represented by an exact
hidden constant $1/2$, propagated using weight
$2\operatorname{arctanh}(1/2)$; its contribution is included in $B$.
At the last threshold, use twice the affine form as the linear readout.
It has signed value at least $3/2$. This construction has exactly $q$
weighted layers, no later biases, and polynomial width and weights.

The quantified fixed-intrinsic-depth theorem then gives a sample bound
$S^{O_\zeta(q)}$ times the displayed polynomial in $L$ and the
statistical parameters. Since $q=O(1+h/\log\log S)$, this is the
claimed estimate. Its small-width part is already expressed using the
actual complexity $C_L(f)$, so the alternative circuit's larger width
does not impose an additional lower-width assumption.
\end{proof}

This is a circuit-size/depth tradeoff rather than a polynomial-product
representation theorem. In particular, for balanced recursive majority
it improves the available quasipolynomial upper bound but does not
settle whether a polynomial bound is possible.

\section{Formula targets independently of their original depth}
\label{sec:formula-balancing}

Formula balancing is classical, following Spira~\cite{spira1971}.
We include the elementary construction needed here to make the resulting
sample bound explicit. This section does not claim a new circuit-balancing
theorem.

\begin{lemma}[Polynomial-size balancing of a Boolean formula]
\label{lem:formula-balancing}
An AND/OR/NOT formula with $s$ gates and literal occurrences has an
equivalent fan-in-two circuit of size polynomial in $s+2$ and depth
$O(\log(s+2))$, regardless of the original formula depth.
\end{lemma}

\begin{proof}
Push negations to the leaves using De Morgan's laws, eliminate double
negations, and simplify constants. This does not increase the number
of literal occurrences. Binarize any larger-fan-in AND/OR gates; this
adds only linearly many gates and does not change the leaf count.
A constant formula is immediate. Otherwise
let $n$ denote its number of literal occurrences, counting repeated
occurrences separately.

For $n\ge2$ there is a subformula $g$ with between $n/3$ and $2n/3$
literal occurrences. To find one, descend from the root while a child
has more than $2n/3$ leaves. At the last such node both children have
at most $2n/3$ leaves, and at least one has at least $n/3$ leaves.
Regard this one occurrence of $g$ as a formal Boolean input. After
substituting either value $b\in\{0,1\}$ and simplifying constants,
the remaining formula $F_b$ has at most $n-|g|\le2n/3$ literal
occurrences. The Shannon identity is
\[
 F=(g\wedge F_1)\vee(\neg g\wedge F_0).
\]
It is valid even if variables appearing in $g$ also occur elsewhere
in the original formula.

Recursively balance $g,F_0,F_1$ and join the three resulting circuits
by the four indicated Boolean gates. Their depths increase by at most
three, and each recursive subproblem has at most $2n/3$ occurrences.
Hence the depth is $O(\log(n+1))$. There are at most three recursive
subproblems per call, so the size is polynomial; for example a bound
$O((n+1)^4)$ follows immediately by induction. Reuse the balanced
output of $g$ in its two appearances in the displayed identity.
This establishes the asserted circuit bounds.
\end{proof}

\begin{corollary}[A quasipolynomial bound for arbitrary formulas]
\label{cor:arbitrary-formula-learning}
Let $f$ have a Boolean formula with at most $s$ gates and literal
occurrences, and put $S=\max\{16,2+d+s\}$. There is a tanh realization
with unit margin, polynomial-in-$S$ width and physical weight bounds,
and depth $q$ satisfying
\[
 2\le q\le C\frac{\log S}{\log\log S}
\]
for a universal constant $C$. For each fixed $\zeta>0$ and every
learner depth $L\ge q$, the unrestricted posterior sample bound is
\[
 \exp\!\left[O_\zeta\!\left(\frac{\log^2S}{\log\log S}\right)\right]
 \bigl(L+\varepsilon^{-1}+\log(2/\delta)\bigr)^{O_\zeta(1)},
\]
simultaneously for all actual widths $N\ge C_L(f)$ and all
$0<\sigma^2\le N^{-(1+\zeta)}$.
This applies in particular to read-once formulas and does not depend
on the depth of the given formula.
\end{corollary}

\begin{proof}
Lemma~\ref{lem:formula-balancing} gives a circuit of polynomial size
and depth $O(\log S)$. Apply
Proposition~\ref{prop:boolean-layer-blocking} to that circuit.
The logarithm of its size parameter is $O(\log S)$ and its double
logarithm is comparable to $\log\log S$. Its blocked tanh depth is
therefore $O(\log S/\log\log S)$, with polynomial width and weights.
The quantified intrinsic-depth bound gives the displayed sample
estimate. All-width coverage is inherited from that bound, not from
identifying the target's actual minimum width with the construction.
\end{proof}

When a supplied formula is already shallower, its original blocked
representation can be used instead. The result above is still
quasipolynomial in formula size; it does not resolve polynomial sample
complexity for all polynomial-size formulas.

\section{Polynomial product norm suffices after logarithmic added depth}
\label{sec:product-targets}

The posterior representation ceiling has a partial converse for a
specified finite-width target representation. Polynomial product norm
alone does not balance the physical tanh layers, but small-prefix
linearization and unequal clone counts balance all retained hidden
layers. Only a polynomial readout cost then needs to be transmitted.

\begin{theorem}[Deep targets with polynomial product norm]
\label{thm:product-target-log-padding}
Suppose a width-$C$, depth-$D_0\ge2$ tanh network has global signed
margin at least one for $f$, and product of layer Frobenius norms at
most $P$. Put $P_0=\max\{1,P\}$ and
\[
 \tau^2=\frac1{16P_0D_0(d+1)^{3/2}},\qquad
 T=\frac{32CP_0^2}{\tau^2}.
\]
For every fixed $\zeta>0$, the original unrestricted Gaussian posterior
has sample complexity polynomial in
\[
 d,C,D_0,P_0,L,\varepsilon^{-1},\log(2/\delta)
\]
at every learner depth
\[
 L\ge D_0+\lceil\log T\rceil,
\]
simultaneously for every $N\ge C_L(f)$ and every
$0<\sigma^2\le N^{-(1+\zeta)}$. The same polynomial conclusion holds
under the alternative depth condition $L\ge2D_0$.
There is no bound on the individual weights of the specified target
representation and no internal robustness restriction.
\end{theorem}

\begin{proof}
Write $M_j$ for the actual Frobenius norm of its $j$th matrix, including
the readout, and $S_j=\prod_{i\le j}M_i$ for the prefix products.
All norms are positive, since the global margin excludes a zero map.
The augmented input norm is $R_0=\sqrt{d+1}$.
At every hidden layer, boundedness of tanh and the suffix Lipschitz
bound give
\[
 1\le |z(x)|\le\sqrt C\,\frac P{S_j},
 \qquad S_j\le P\sqrt C.
\]
Using an upper bound $P$ instead of the exact full product only weakens
this inequality.

Replace the tanh activation by the identity at each hidden layer whose
original prefix product satisfies $S_j<\tau$.
The inequality $|u-\tanh u|\le |u|^3/3$ and a telescoping argument give
total output change at most
\[
 \sum_{j:S_j<\tau}\frac{P}{S_j}
      \frac{(S_jR_0)^3}{3}
 \le\frac{PD_0\tau^2R_0^3}{3}\le\frac1{48}.
\]
The same prefix and suffix bounds apply throughout this telescoping,
since both identity and tanh are zero-preserving and one-Lipschitz.
Collapse the resulting consecutive linear matrices into their next
nonlinear layer or the final readout, then double the readout. The
result still computes $f$ with unit margin.

If every hidden activation was removed, $f$ is a Boolean halfspace,
and Theorem~\ref{thm:halfspace-unrestricted-all-depths} proves the claim.
Otherwise let $H\ge1$ be the number of retained hidden layers, with
original indices $j_1<\cdots<j_H$. Their retained prefix bounds satisfy
\[
 \tau\le S_{j_s}\le P_0\sqrt C.
\]
The first collapsed matrix has Frobenius norm at most $S_{j_1}$;
the $s$th has norm at most $S_{j_s}/S_{j_{s-1}}$; and the doubled
readout has norm at most $2P_0/S_{j_H}$. These are assigned norm bounds;
matrix multiplication can make the actual norms smaller without
affecting any subsequent upper bound. In particular every collapsed
physical weight is at most
\[
 R=\max\{1,2P_0\sqrt C/\tau\}.
\]

Define
\[
 a=\frac{d\tau^2}{8C},\qquad
 r_s=\frac{a}{dS_{j_s}^2}=\frac{\tau^2}{8C S_{j_s}^2},
 \qquad r_{\min}=\frac{\tau^2}{8C^2P_0^2}.
\]
Then $r_{\min}\le r_s\le1/(8C)$ and $a<1/8$.
At common width $N$, give every original neuron in retained layer $s$
exactly $c_s=\lceil Nr_s\rceil$ clones, divide outgoing weights by
$c_s$, and fill the remaining positions with zero neurons.
For
\[
 N\ge M:=\left\lceil\max\{2C,8D_0/r_{\min},2R,4T,4\}\right\rceil,
\]
the clone counts fit in width $N$ and
$Nr_s\le c_s\le Nr_s(1+1/(8D_0))$.
This preserves the collapsed network exactly. Its pre-readout energy
and readout energy obey
\[
 E_{\rm pre}\le a\left(1+\frac1{8D_0}\right)
 +(H-1)\left(1+\frac1{8D_0}\right)\le H,
\]
\[
 c^2=N\|v_N\|_2^2
 \le\frac{dS_{j_H}^2}{a}\frac{4P_0^2}{S_{j_H}^2}
 =\frac{32CP_0^2}{\tau^2}=T.
\]
For the first inequality, $H\le D_0-1$ and $a<1/8$ leave more than
enough slack to absorb all clone rounding.

Replace the readout by the scalar-transmission construction with
$m=L-H$ weighted layers, and double its final readout as in the
shallow-target proof. The resulting unit-margin witness has energy
\[
 K\le H+(m+15)T^{1/m}.
\]
The stated depth conditions ensure $m\ge2$.
If $L\ge D_0+\lceil\log T\rceil$, then $m\ge\log T$, and
$e^u\le1+2u$ for $0\le u\le1$ gives
\[
 K\le L+15+32\log T.
\]
Hence the Gaussian-distance parameter is bounded by
\[
 Q_*^2=((K+2)/L)^L\le e^{17}T^{32}.
\]
Under the alternative condition $L\ge2D_0$, one has $L/m\le2$ and
the direct estimate
\[
 K+2\le(L+17)T^{1/m},\qquad Q_*^2\le e^{17}T^2.
\]
In either case, $K,Q_*$ and $M$ are polynomial in the stated target
parameters and learner depth. The Gaussian-distance theorem therefore
gives the required sample bound for all $N\ge M$.

We check the actual smaller-width convention. The collapsed weights
are at most $R\le M/2$, and cloning can only reduce internal and
readout weights. The first scalar-transmission matrix has weights at
most $1.1R<M$. If $q^2=\max\{c^2,1\}^{1/m}\le T$, the remaining
scalar weights have magnitudes at most $2.2q/N$ and $2q/N$.
Since $M\ge4T\ge4$, these are less than $M$.
Thus the constructed width-$M$ target is admissible and $C_L(f)\le M$.
For $C_L(f)\le N<M$, clone a defining minimum-width target and double
its readout; its width and weights are bounded by a polynomial in $M$.
The finite-width theorem supplies the same polynomial sample form.
Intersect the two sample events, with failure probabilities $\delta/2$
each, to conclude.
\end{proof}

The theorem is polynomial in $P$ itself. For a generic bounded-weight
deep circuit, $P$ may still grow exponentially with its depth. Thus this
does not settle the unrestricted general small-circuit conjecture.
It does show that a polynomial product-norm representation is a
sufficient target condition after a logarithmic number of additional
layers, without imposing that condition on the support of the prior.

\section{Width compression with controlled product norm}
\label{sec:pf-width-compression}

The width dependence in Theorem~\ref{thm:product-target-log-padding}
can be removed. Boundedness and near-linearity of tanh permit a
norm-preserving compression: very large normalized rows can be dropped,
very small rows can be linearized, and the remaining neurons can be
sampled with a probability chosen to balance their two adjacent norms.

\begin{theorem}[Finite-set width compression for tanh]
\label{thm:pf-width-compression}
Let $X$ be a nonempty finite set of inputs of norm at most $R\ge1$.
Suppose a depth-$D_0\ge2$ tanh network with linear readout, arbitrary
finite hidden widths, and no later biases has product of Frobenius norms
at most $P_0\ge1$. For every $0<\epsilon<1$, it has a same-depth
approximation $\widetilde z$ such that
\[
 \max_{x\in X}|z(x)-\widetilde z(x)|\le\epsilon,
 \qquad \prod_{j=1}^{D_0}\|\widetilde W^j\|_F\le2P_0,
\]
and every new hidden width is bounded by a universal polynomial in
\[
 P_0,D_0,R,\log(2|X|),\epsilon^{-1}.
\]
The bound is independent of the original widths and individual weights.
An augmented constant input is allowed and counts toward $R$.
\end{theorem}

\begin{proof}
Put $H=D_0-1$, $P_*=2P_0$, $e=\epsilon/(8H)$ and
$\eta=1/(4H)$. We compress the hidden layers in order from the first
to the last. At each step the product norm increases by at most
$1+\eta$, and the output changes by at most $6e$ on $X$.
Thus every intermediate product is below $P_*$, and the final product
is at most $e^{1/4}P_0<2P_0$, while total error is less than $\epsilon$.
A zero matrix makes the entire network zero and is trivial, so all
prefix and pair norms used below may be assumed positive. If a
compression step produces a zero matrix, stop and use the zero network
with width one; the accumulated approximation estimates still apply.

Consider a hidden matrix $W$ and its immediately following matrix $V$;
the latter may be the final linear readout. Let $S$ be the product of
all Frobenius norms strictly before $W$, with $S=1$ at the first layer,
and put
\[
 A=\|W\|_F\|V\|_F.
\]
The incoming feature vector $h(x)$ has norm at most $SR$. The part
strictly after $V$ has Lipschitz constant at most $P_*/(SA)$.
If $w_i^T$ is row $i$ of $W$ and $v_i$ is column $i$ of $V$, write
$a_i=\|w_i\|_2$ and $b_i=\|v_i\|_2$.
Zero rows or zero outgoing columns can be deleted.

Set
\[
 B=P_*/e,\qquad \tau^2=3e/(P_*R^3),\qquad
 \Delta=e/(P_*R),\qquad Q=(B/\tau)^2.
\]
Partition neurons into large, small and middle according as
$Sa_i>B$, $Sa_i<\tau$, or $\tau\le Sa_i\le B$.
There are at most $S^2\|W\|_F^2/B^2$ large neurons. Their activation
vector has norm at most the square root of this count. Deleting them
therefore changes the final output by at most $P_*/B=e$.

For the small neurons, replace tanh by the identity. Since
$|u-\tanh u|\le |u|^3/3$, their combined activation error has norm at
most $\tau^2R^3S\|W\|_F/3$. Propagation through the suffix bounds
the output error by $P_*\tau^2R^3/3=e$.
Their resulting contribution to the next preactivation is the linear map
\[
 B_0=\sum_{i\ \mathrm{small}}v_iw_i^T,
 \qquad \|B_0\|_*\le\sum_{i\ \mathrm{small}}a_ib_i,
\]
where $\|\cdot\|_*$ is the nuclear norm.

Let $A_m=\sum_{i\ \mathrm{middle}}a_ib_i$.
If $A_m<\Delta A$, delete all middle neurons: their output contribution
is at most $P_*R A_m/A\le e$. Otherwise $A_m\ge\Delta A>0$, and
we approximate their contribution by sampling. Define
\[
 Z=\sum_{i\ \mathrm{middle}}b_i/a_i,
 \qquad p_i=(b_i/a_i)/Z.
\]
Draw $m$ independent indices from $(p_i)$. A sampled neuron keeps its
incoming row $w_i^T$ and has outgoing column $v_i/(mp_i)$.
If $\widehat A=(Z/m)\sum_{r=1}^m a_{i_r}^2$ and $c=m/Z$, its new
incoming and outgoing squared Frobenius norms are exactly
\[
 c\widehat A\quad\hbox{and}\quad \widehat A/c.
\]
Moreover $\mathbb E\widehat A=A_m$. Since middle row norms lie between
$\tau/S$ and $B/S$, one has
\[
 Z\max_i a_i^2\le A_m Q\le AQ.
\]
The random vector $v_i\tanh(w_i^Th(x))/p_i$ has norm at most
$AQSR$, uniformly on $X$.

Consequently some choice of the samples has both
\[
 \widehat A\le A_m+\eta A
\]
and middle-preactivation approximation error at most $eSA/P_*$
simultaneously for every $x\in X$, provided
\[
 m\ge64Q^2\max\{\eta^{-2},(P_*R/e)^2\}
                       \log(8(|X|+1)).
\]
Here and below take integer ceilings. For detail, the scalar random
variables $Za_i^2/A$ are bounded by $Q$, so Hoeffding controls their
mean. For a Hilbert-valued random vector bounded by $AQSR$, the mean
deviation has expectation at most $AQSR/\sqrt m$; bounded differences
gives a dimension-free tail. A union bound over $X$ and the scalar
condition has positive success probability with the displayed constant.
Propagation of the vector error gives final output error at most $e$.

We next compress the linear branch. Truncate the singular-value
decomposition of $B_0$ to rank
\[
 k=\lceil P_*R/e\rceil,
\]
obtaining $B_k=UV_0$ with
$\|U\|_F^2=\|V_0\|_F^2=\nu:=\|B_k\|_*\le\|B_0\|_*$.
The residual operator norm is at most $\|B_0\|_* /(k+1)\le A/(k+1)$,
so its final output error is at most $P_*R/(k+1)\le e$.
If $\nu=0$, omit the linear branch entirely.

When the middle branch was sampled, retain $c=m/Z$. When it was
deleted and $\nu>0$, choose $c=1/(\nu S^2)$ instead.
Represent $B_k$ approximately using $K$ copies of each of its at most
$k$ singular channels: their incoming rows are the rows of
$\sqrt{c/K}\,V_0$, repeated $K$ times, and their outgoing columns are
the corresponding columns of $U/\sqrt{cK}$.
Their total incoming and outgoing squared norms are $c\nu$ and
$\nu/c$. Their linear limit is exactly $B_kh$.
The tanh linearization error at the next preactivation has norm at most
\[
 \frac{c\nu^2(SR)^3}{3K}.
\]
After propagation this is at most
$P_*c\nu S^2R^3/(3K)$, using $\nu\le A$.
In the sampled case,
\[
 Z\ge A_mS^2/B^2,
 \qquad c\nu S^2\le mB^2/\Delta;
\]
in the deleted case this quantity equals one. Hence it suffices to use
\[
 K=\left\lceil\frac{P_*R^3}{e}
                   \max\{1,mB^2/\Delta\}\right\rceil.
\]
The linear-channel approximation then contributes at most $e$ error.
When there is no sampled middle branch, the same bound is harmless;
one may use the displayed deterministic value of $m$ merely to define
a common width upper bound.

Combining the sampled and linear branches gives pair product norm
\[
 \widehat A+\nu
 \le A_m+\sum_{i\ \mathrm{small}}a_ib_i+\eta A
 \le(1+\eta)A.
\]
The last inequality is Cauchy--Schwarz over the original rows and
columns. If the middle branch was deleted, the pair product is simply
$\nu\le A$. These are exact norm calculations for the constructed
tanh matrices, not only bounds on their linear limits.

The compressed hidden width is at most $m+kK$. All quantities in this
bound are universal polynomials in $P_0,H,R,e^{-1}$ and
$\log(2|X|)$. Crucially, no original dimension of either adjacent hidden
layer occurs. Only the two adjacent matrices are altered. Processing
the next hidden layer changes its own width without undoing the width
bound for the preceding one. Uniform zero-padding at the end makes
all hidden widths equal if desired. Summing the at most six errors and
multiplying the pair-product bounds proves the theorem.
\end{proof}

\begin{corollary}[Uniform compression on a bounded input set]
\label{cor:pf-compact-compression}
If the input dimension is $q$ and $X$ is any subset of the Euclidean
ball of radius $R\ge1$, the same conclusion holds with width polynomial
in $q,P_0,D_0,R,\epsilon^{-1}$, independently of whether $X$ is finite.
\end{corollary}

\begin{proof}
Choose an $\epsilon/(6P_0)$-net of the radius-$R$ ball with at most
$(1+12P_0R/\epsilon)^q$ points. Apply the finite-set theorem to this net
with error $\epsilon/2$. The original and compressed networks have
Lipschitz constants at most $P_0$ and $2P_0$, respectively. Their
difference therefore changes by at most $\epsilon/2$ between any point
and its nearest net point. This proves uniform error at most $\epsilon$.
The logarithm of the net size is polynomial in the displayed quantities.
\end{proof}

\begin{corollary}[A width-free sufficient representation condition]
\label{cor:width-free-pf-sufficiency}
Suppose $f$ has some finite-width, depth-$D_0$ tanh representation with
unit global margin and product of Frobenius norms at most $P_0\ge1$.
For each fixed $\zeta>0$, there are universal polynomials and a universal
constant $c$ such that its unrestricted Gaussian posterior sample
complexity is polynomial in
\[
 d,D_0,P_0,L,\varepsilon^{-1},\log(2/\delta)
\]
whenever
\[
 L\ge D_0+c\log(2+d+D_0+P_0),
\]
simultaneously for every $N\ge C_L(f)$ and every
$0<\sigma^2\le N^{-(1+\zeta)}$.
The same conclusion holds at every $L\ge2D_0$.
There is no dependence on the original width or on individual weights.
\end{corollary}

\begin{proof}
Apply Theorem~\ref{thm:pf-width-compression} with the augmented Boolean
cube, $R=\sqrt{d+1}$, $|X|=2^d$, and approximation error $1/4$.
The resulting representation has width polynomial in $d,D_0,P_0$,
the same depth, and global signed margin at least $3/4$.
Double its readout to obtain unit margin and product norm at most
$4P_0$. Theorem~\ref{thm:product-target-log-padding} then applies.
Its $T$ is polynomial in $d,D_0,P_0$, so the required logarithmic
overhead has the displayed form after increasing a universal constant.
\end{proof}

Together with the posterior ceiling, this identifies polynomial product
norm as a substantive representation criterion for this prior. The
necessity statement concerns constant-error approximations at the same
learner depth; the sufficiency statement above assumes an exact Boolean
representation and permits logarithmically many added layers. These
qualifications prevent treating the two statements as an exact
equivalence.

\section{What width compression does and does not say about circuits}
\label{sec:product-normal-form}

\begin{corollary}[A finite-description normal form]
\label{cor:product-finite-description}
Suppose a Boolean function on $\{-1,1\}^d$ has an arbitrary finite-width,
depth-$D\ge2$ tanh representation with global signed margin at least one
and product of Frobenius norms at most $P\ge1$. There is a representation
of the same function, with margin at least one, using at most $D$
weighted layers, polynomial width and polynomially bounded individual
weights in $d,D,P$, and product norm at most $8P$.
Here a purely affine function of the augmented input is allowed as the
case of one weighted layer. The weights may moreover be chosen dyadic
with polynomial bit length, at the expense of replacing the product
bound by $16P$. Consequently the function has a Boolean circuit of size
polynomial in $d,D,P$, with no dependence on the original widths,
weights, or their precision.
\end{corollary}

\begin{proof}
Apply Theorem~\ref{thm:pf-width-compression} on the augmented cube with
error $1/4$, and double its readout. This gives width
$C=\operatorname{poly}(d,D,P)$, product at most $P_1=4P$, and signed
margin at least $3/2$. Apply the small-prefix linearization from the
proof of Theorem~\ref{thm:product-target-log-padding}, with
\[
 \tau^2=\frac1{16P_1D(d+1)^{3/2}}.
\]
It changes every output by at most $1/48$. Collapse the resulting linear
blocks and double the final readout. The resulting signed margin exceeds
two, the product is at most $2P_1=8P$, and every physical weight is at
most
\[
 R=\max\{1,2P_1\sqrt C/\tau\}
   =\operatorname{poly}(d,D,P).
\]
The number of weighted layers does not increase. If every hidden
activation is removed, the remaining affine map has norm at most
$2P_1$ and satisfies the same claims.

For completeness, polynomial precision can be imposed without assuming
that the displayed real weights were efficiently computable. Let $D'$
be the remaining depth and choose a polynomial $B\ge2$ that bounds the
Frobenius norm and the row absolute sum of every remaining matrix, and
also bounds all input and hidden dimensions. Each matrix norm $M_j$
satisfies
\[
 M_j\ge\frac1{\sqrt{d+1}\,B^{D'-1}}.
\]
Indeed the margin gives
$1\le |z(x)|\le\sqrt{d+1}\prod_iM_i$, and all other norms are at most
$B$. Round each entry to a dyadic mesh $2^{-b}$, where
\[
 b=O\bigl(D'\log B+\log(d+D'+2)\bigr)
\]
with a sufficiently large universal constant. The matrix Frobenius
perturbations are then at most $M_j/(4D')$, so the product increases
by at most $(1+1/(4D'))^{D'}<2$. The standard layerwise error recursion,
using the row-sum bound $B$ and the one-Lipschitz property of tanh,
makes the total output change less than $1/2$ at the same precision.
The initial margin exceeding two therefore leaves margin at least one.
There are polynomially many entries and $b$ is polynomial in $d,D,P$.

Finally, evaluate each affine map and tanh to sufficiently many binary
digits and propagate a total error smaller than $1/2$. Arithmetic on
these polynomial-length numbers and approximation of tanh on the
polynomially bounded preactivation range use polynomial-size Boolean
circuits. The output sign remains exact because of the margin.
\end{proof}

The converse direction is not supplied by width compression: a
polynomial-size Boolean circuit is not thereby given a polynomial
product-norm tanh representation. In particular, widening an existing
implementation is not itself a proof of such a bound. If a polynomial
product-norm implementation exists at any finite width, the corollary
shows that a polynomial-width, finite-description implementation exists
as well.

Nor does this normal form impose constant threshold depth. The
rank-one approximation argument gives retained depth of order
$P^4\log(2P)$, after restriction to the at most $P^2$ relevant
coordinates, which is only a polynomial depth bound when $P$ is
polynomial. Its combination with width compression therefore gives
no inclusion in $\mathrm{TC}^0$ or $\mathrm{NC}^1$ and no separation of
those circuit classes.

\section{Product norm controls Boolean description size independently of depth}
\label{sec:boolean-product-depth-saturation}

The rank-one bottleneck strategy of Golowich--Rakhlin--Shamir
\cite{golowich2018} underlies depth-independent norm-based complexity
bounds. Here we use it for a more specific representation statement:
global Boolean margin permits a short, polynomial-norm tanh
implementation of the scalar suffix, without carrying a constant rail
through the retained prefix. The conclusion bounds a finite Boolean
description independently of the original depth as well as width.

\begin{lemma}[Relevant coordinates and freezing]
\label{lem:pf-junta-freezing}
A Boolean function with a globally unit-margin tanh representation
of product norm at most $P$ depends on at most $P^2$ coordinates.
After deleting all irrelevant coordinates, it has a representation
of the same depth and no larger product norm.
\end{lemma}

\begin{proof}
Let $T=\prod_{j\ge2}\|W^j\|_F$. For every relevant coordinate $i$,
some cube edge in that coordinate has opposite labels at its endpoints.
The two outputs differ by at least two, whereas the Lipschitz bound
along that edge is $2\|W^1e_i\|_2T$. Thus
$\|W^1e_i\|_2T\ge1$. Squaring and summing over the relevant coordinates
gives their number at most $\|W^1\|_F^2T^2\le P^2$.

Freeze each irrelevant coordinate $i$ to an independent uniform sign
$a_i$. This replaces the first-layer bias column $w_0$ by
$w_0+\sum_i a_iw_i$, while retaining the relevant columns. Its expected
squared Frobenius norm is exactly the original squared Frobenius norm.
Some assignment therefore does not increase that norm. Since the
function is independent of these coordinates, the reduced network
retains its unit margin on the full reduced cube.
\end{proof}

\begin{lemma}[A bias-free odd scalar classifier]
\label{lem:odd-scalar-classifier}
Let $P,R\ge1$ and let $g:[-R,R]\to\mathbb R$ be odd and $P$-Lipschitz.
Let $T\subset[-R,R]$ be finite and nonempty, with
$|g(t)|\ge7/8$ for $t\in T$. There is a tanh network with two hidden
layers, a linear readout and no biases, whose sign agrees with $g$
on $T$ and whose signed output there is at least $7/4$.
Its width, weights, and product of Frobenius norms are bounded by
universal polynomials in $P,R$.
\end{lemma}

\begin{proof}
Extend the set to $T\cup(-T)$ and sort its positive absolute values.
All are at least $7/(8P)$, since $g(0)=0$. Consecutive values with
opposite signs of $g$ are separated by at least $7/(4P)$.
Place a threshold $\theta_j$ at the midpoint of each such transition
gap. There are $m\le2PR$ transitions, every $0<\theta_j\le R$,
and every actual absolute value is at distance at least $7/(8P)$
from every threshold.

The resulting odd step classifier has the exact expression
\[
 F(t)=c_0\operatorname{sign}(t)+
   \sum_{j=1}^m c_j\operatorname{sign}
       \bigl(t-\theta_j\operatorname{sign}(t)\bigr),
 \qquad |c_0|\le1,\quad c_j\in\{-1,1\}.
\]
Indeed, the positive-side label changes by $2c_j$ at transition $j$.
The changes alternate, so the coefficient
$c_0=\ell_0+\sum_jc_j$ is in $\{-1,0,1\}$, where $\ell_0$ is the
first positive-side label. The displayed formula is the odd extension
of this step classifier and equals $\operatorname{sign}g(t)$ on $T$.

Set
\[
 \rho=\frac1{2P},\qquad
 \delta=\min\left\{\frac\rho{4R},\frac1{32(m+1)}\right\},\qquad
 A=\frac1{2\rho}\log\frac2\delta,
\]
\[
 \alpha^2=\frac\rho{4R^3},\qquad
 B=\rho^{-1}\log(64(m+1)),\qquad
 c=\operatorname{arctanh}(1/2).
\]
The first hidden layer computes
$u(t)=\tanh(\alpha t)$ and $s(t)=\tanh(At)$.
On $T\cup(-T)$,
\[
 |u(t)/\alpha-t|\le\rho/12,
 \qquad |s(t)-\operatorname{sign}(t)|\le\delta.
\]
Thus replacing $t-\theta_j\operatorname{sign}(t)$ by
$u(t)/\alpha-\theta_js(t)$ changes it by at most
$\rho/12+R\delta\le\rho/3$. The original magnitude exceeds $\rho$,
so its sign is unchanged and its new magnitude exceeds $\rho/2$.
The second hidden layer computes
\[
 r(t)=\tanh(cs(t)),\qquad
 v_j(t)=\tanh\bigl(B[u(t)/\alpha-\theta_js(t)]\bigr).
\]
Consequently $2r(t)$ approximates $\operatorname{sign}(t)$ to error
at most $2\delta$, and each $v_j$ approximates its desired sign to
error at most $1/[32(m+1)]$. Therefore
\[
 z(t)=2\left[2c_0r(t)+\sum_{j=1}^m c_jv_j(t)\right]
\]
has signed value at least
$2[1-2\delta-m/(32(m+1))]\ge15/8>7/4$ on $T$.
This includes $m=0$, when only the sign rail is needed.

The three squared matrix norms are bounded respectively by
\[
 \alpha^2+A^2,\qquad
 c^2+mB^2(\alpha^{-2}+R^2),\qquad 16+4m.
\]
Since $m\le2PR$, these expressions, their product, and every physical
weight are polynomially bounded in $P,R$. All biases are zero.
\end{proof}

\begin{theorem}[A depth-independent Boolean normal form]
\label{thm:depth-independent-pf-normal-form}
There is a universal polynomial $p$ with the following property.
Suppose $f:\{-1,1\}^d\to\{-1,1\}$ has some finite-depth,
finite-width tanh representation with unit global margin and product
of Frobenius norms at most $P\ge1$. Then $f$ has a globally
unit-margin representation using only its relevant input coordinates,
with depth, hidden widths, individual weights, and product of
Frobenius norms at most $p(P)$. Its weights can be chosen dyadic
with at most $p(P)$ bits each. The bounds are independent of the
original depth, widths, individual weights and their precision.
There are at most $P^2$ relevant coordinates.
\end{theorem}

\begin{proof}
First use Lemma~\ref{lem:pf-junta-freezing}. If no coordinate is
relevant, a constant network of width one and depth two suffices.
Otherwise write $k\le P^2$ for the number of remaining coordinates
and $R=\sqrt{k+1}$. Let $D$ be the remaining original depth, and put
\[
 r=\left\lceil128P^2R^2\log(2PR)\right\rceil.
\]
If $D\le r+2$, retain the network for the moment. Otherwise we
produce a representation with at most $r+2$ weighted layers and
polynomial product norm.

Write $M_j=\|W^j\|_F$, $s_j=\|W^j\|_{\rm op}$ and
$\operatorname{sr}_j=M_j^2/s_j^2$. Unit margin implies
$\prod_js_j\ge1/R$, because all activations are zero-preserving
and one-Lipschitz. Hence
\[
 \sum_j\log\operatorname{sr}_j
 =2\log\frac{\prod_jM_j}{\prod_js_j}
 \le2\log(PR).
\]
Among the first $r$ hidden layers choose one, $j$, with
$\log\operatorname{sr}_j\le2\log(PR)/r$.
Replace $W^j$ by its leading singular component $s_juv^T$.
The discarded operator norm is at most
$M_j\sqrt{1-1/\operatorname{sr}_j}
 \le M_j\sqrt{\log\operatorname{sr}_j}$.
Propagation through the original prefix and suffix therefore changes
every output by at most
\[
 PR\sqrt{2\log(PR)/r}\le1/8.
\]
The modified network retains the target's signs and margin at least
$7/8$.

Put $S=\prod_{i<j}M_i$, with $S=1$ if $j=1$, and normalize the
bottleneck scalar as $t(x)=v^Th^{j-1}(x)/S$. Its absolute value is at
most $R$. Everything from the rank-one matrix onward is an odd
scalar function $g(t)$ with Lipschitz constant at most
$s_jS\prod_{i>j}M_i\le P$. Apply
Lemma~\ref{lem:odd-scalar-classifier} on the finite set of these
actual scalar values.

Attach its two-hidden-layer classifier to the original prefix.
The first two new rows are $\alpha v^T/S$ and $Av^T/S$.
Their common factor $1/S$ cancels the prefix product $S$ in the
total Frobenius product. Thus the complete new network has polynomial
product norm in $P,R$, independent of every original dimension and
weight scale. Its depth is $j+2\le r+2$, and it has the required
global signs with margin at least $7/4$.

In either case we now have depth and product norm polynomial in $P$,
while widths may still be arbitrary. Apply the finite-description
normal form, Corollary~\ref{cor:product-finite-description}, with
input dimension $k\le P^2$. It gives polynomial width, physical
weights and dyadic bit length, keeping unit margin and a polynomial
product norm. Restore zero columns for the irrelevant original
coordinates if a $d$-input network is desired. If the normal form is
purely affine, apply tanh to that affine output and append the
readout $2$; its unit margin is preserved and all bounds remain
polynomial, now within the usual architecture with a hidden layer.
\end{proof}

\begin{corollary}[Counting globally margined product-norm classifiers]
\label{cor:all-depth-pf-counting}
Let $\mathcal H_P$ contain all Boolean functions admitting any
finite-depth, finite-width globally unit-margin tanh representation
with product norm at most $P\ge1$. For a universal polynomial $p$,
\[
 \log|\mathcal H_P|\le p(P)+P^2\log(d+1).
\]
Consequently, with probability at least $1-\delta$, every member
of $\mathcal H_P$ consistent with $n$ independent realizable examples
has error at most $\varepsilon$, provided
\[
 n\ge\frac{p(P)+P^2\log(d+1)+\log(1/\delta)}{\varepsilon}.
\]
There is no depth or width restriction in this simultaneous event.
\end{corollary}

\begin{proof}
There are at most $(d+1)^{\lfloor P^2\rfloor}$ possible ordered
lists of relevant coordinates, allowing a padding symbol. Once the
coordinates are chosen, Theorem~\ref{thm:depth-independent-pf-normal-form}
encodes the entire polynomial-size dyadic network in $p(P)$ bits,
after enlarging a universal polynomial. Count these descriptions and
apply the usual union bound: an error-$\varepsilon$ classifier is
consistent with $n$ independent examples with probability at most
$e^{-n\varepsilon}$.
\end{proof}

In particular, counting implies that most Boolean functions of $k$
variables require product norm at least $2^{c k}$ for some universal
$c>0$, even when arbitrary depth and width are allowed. This is not
a small-circuit lower bound: those generic functions also require
exponentially large ordinary Boolean descriptions.

\section{A shallow product-norm obstruction that depth removes}
\label{sec:shallow-product-separation}

One possible route to a stronger normal form would replace every
small-product deep network by a one-hidden-layer network with polynomial
product cost. This is false, even for a Boolean target that the present
unrestricted posterior learns with polynomial sample complexity.
The obstruction concerns shallow representation, rather than the
learner's ability to use a different deep representation.

For $x,y\in\{-1,1\}^n$, put $u_i=(1-x_i)/2$, $v_i=(1-y_i)/2$, and define
\[
 \mathrm{IP}_n(x,y)=(-1)^{\sum_{i=1}^n u_iv_i}.
\]
This is the Boolean inner-product function on $2n$ input coordinates.
As elsewhere, depth counts weighted layers: one hidden layer has depth
two, and two hidden layers have depth three.

\begin{theorem}[Exponential shallow cost, polynomial deep cost]
\label{thm:ip-shallow-product-separation}
For every $n\ge1$, every finite-width one-hidden-layer tanh network
\[
 z(x,y)=\sum_{a=1}^m v_a\tanh\bigl(w_a^T(1,x,y)\bigr)
\]
with $\mathrm{IP}_n(x,y)z(x,y)\ge1$ on the entire cube satisfies
\[
 \|v\|_2\|W\|_F
 \ge \frac{2^{n/2}}{(2n+3)\sqrt{2n+1}},
 \qquad W=(w_a^T)_{a=1}^m.
\]
There is, however, a depth-three, width-$(n+1)$ tanh realization in the
present architecture, with signed margin at least one and
\[
 \prod_{\ell=1}^3\|W^\ell\|_F
 \le20n^{5/2}\log^2(8n).
\]
It uses only the augmented first-layer input, with no later biases.
Its individual weights are polynomially bounded in $n$.
\end{theorem}

We first prove a correlation estimate. All expectations in this section
are under independent uniform Boolean inputs.

\begin{lemma}[Halfspaces against a Hadamard sign matrix]
\label{lem:ip-halfspace-correlation}
Let $N=2^n$ and $H_{uv}=(-1)^{u\cdot v}$ for
$u,v\in\{0,1\}^n$. If $H'$ is obtained from $H$ by independently
changing row and column signs, then every matrix
\[
 A_{uv}=\operatorname{sign}(\alpha_u+\beta_v+c)
\]
satisfies
\[
 \left|\frac1{N^2}\sum_{u,v}H'_{uv}A_{uv}\right|
 \le \kappa_n:=\frac{2n+3}{\sqrt N}.
\]
The choice of sign at zero is immaterial.
\end{lemma}

\begin{proof}
The Hadamard identity $HH^T=NI$ gives
$\|H'\|_{\mathrm{op}}=\sqrt N$. Order the columns by decreasing
$\beta_v$. The positive entries of each row of $A$ then form a prefix.
Thus $A=2B-J$, where $B$ is the indicator of these row prefixes and
$J$ is the all-ones matrix.

Decompose each prefix into aligned dyadic intervals, using at most one
interval of each length $2^\ell$, $0\le\ell\le n$. At a fixed scale,
the resulting part of $B$ is a sum of rectangles
$\mathbf1_{R_I}\mathbf1_I^T$, where the column intervals $I$ are
disjoint and each row belongs to at most one $R_I$. Its nuclear norm
is at most
\[
 \sum_I\sqrt{|R_I|2^\ell}
 \le\sqrt{\frac{N}{2^\ell}\,2^\ell\sum_I|R_I|}
 \le N.
\]
Summing over the $n+1$ scales and using $\|J\|_*=N$ gives
\[
 \|A\|_*\le2\|B\|_*+\|J\|_*\le N(2n+3).
\]
By nuclear/operator norm duality, the normalized inner product is at
most $N^{-2}\|H'\|_{\mathrm{op}}\|A\|_*\le\kappa_n$.
Column permutations preserve the norms used in this argument.
\end{proof}

\begin{proof}[Proof of the lower bound in
Theorem~\ref{thm:ip-shallow-product-separation}]
Every affine threshold of $(1,x,y)$ has the matrix form in
Lemma~\ref{lem:ip-halfspace-correlation}. The function
$\phi(t)=\operatorname{sech}^2(t)$ tends to zero at both infinities
and has total variation two. Hence
\[
 \phi(s)=\frac12\int_{\mathbb R}\phi'(t)
                   \operatorname{sign}(s-t)\,dt,
 \qquad \frac12\int_{\mathbb R}|\phi'(t)|\,dt=1.
\]
Integrating the halfspace correlation estimate therefore shows that
\[
 \left|\mathbb E H'\operatorname{sech}^2
                   \bigl(w^T(1,x,y)\bigr)\right|\le\kappa_n
\]
for every $w$ and every allowed row/column sign change $H'$.

Write $\bar x=(1,x,y)$. Since tanh vanishes at zero,
\[
 \tanh(w^T\bar x)
 =\int_0^1\sum_{i=0}^{2n}w_i\bar x_i
                  \operatorname{sech}^2(sw^T\bar x)\,ds.
\]
Multiplying $H$ by any coordinate $\bar x_i$ changes only row signs,
only column signs, or neither. Consequently
\begin{equation}\label{eq:ip-tanh-weight-correlation}
 \left|\mathbb E\mathrm{IP}_n\tanh(w^T\bar x)\right|
 \le\kappa_n\|w\|_1
 \le\kappa_n\sqrt{2n+1}\|w\|_2.
\end{equation}
The unit-margin hypothesis gives $\mathbb E\mathrm{IP}_n z\ge1$.
Sum~\eqref{eq:ip-tanh-weight-correlation} over the hidden neurons and
apply Cauchy--Schwarz:
\[
 1\le\kappa_n\sqrt{2n+1}\sum_a|v_a|\|w_a\|_2
 \le\kappa_n\sqrt{2n+1}\|v\|_2\|W\|_F.
\]
This proves the lower bound, with no restriction on hidden width or
individual coefficient magnitudes. It also lower-bounds the shallow
variation cost $\sum_a|v_a|\|w_a\|_2$ by the same quantity.
\end{proof}

\begin{proof}[Proof of the upper bound in
Theorem~\ref{thm:ip-shallow-product-separation}]
Set $\lambda=\log(8n)$ and $c=\operatorname{atanh}(1/2)$.
The first hidden layer computes
\[
 h_i=\tanh\bigl(\lambda(-x_i-y_i-1)\bigr),\quad 1\le i\le n,
 \qquad h_0=\tanh(c)=\tfrac12.
\]
Let $s_i=+1$ exactly when $u_i=v_i=1$. The corresponding affine
preactivation has absolute value at least one, so
\[
 |h_i-s_i|\le2e^{-2\lambda}=\frac1{32n^2}.
\]
The second hidden layer computes
\[
 r_j=\tanh\left(\lambda\left[\sum_{i=1}^n h_i+n-2j+1\right]\right),
 \quad 1\le j\le n,
 \qquad r_0=\tanh(2ch_0)=\tfrac12.
\]
The constant term in row $j$ is supplied by weight
$2\lambda(n-2j+1)$ from $h_0$. Put $m=\sum_i u_iv_i$. The intended
second-layer sign is $q_j=\operatorname{sign}(2m-2j+1)$.
The perturbation of this odd integer is at most $1/(32n)$, so
its magnitude remains at least $3/4$. It follows that
\[
 |r_j-q_j|\le2e^{-3\lambda/2}\le\frac1{8n}.
\]
With $b_0=1+\sum_{j=1}^n(-1)^j\in\{0,1\}$, the identity
\[
 (-1)^m=b_0+\sum_{j=1}^n(-1)^jq_j
\]
shows that the linear readout
\[
 z=2\sum_{j=1}^n(-1)^jr_j+4b_0r_0
\]
has signed margin at least $2(1-1/8)=7/4$.

The first squared Frobenius norm is $3n\lambda^2+c^2\le4n\lambda^2$.
For the second it is
\[
 n^2\lambda^2+4\lambda^2\sum_{j=1}^n(n-2j+1)^2+4c^2
 =n^2\lambda^2+\frac{4n(n^2-1)}3\lambda^2+4c^2
 \le4n^3\lambda^2.
\]
The readout squared norm is $4n+16b_0^2\le20n$. Thus the three
norms are bounded by $2\sqrt n\lambda$, $2n^{3/2}\lambda$, and
$5\sqrt n$, whose product is the asserted upper bound. The
construction has uniform hidden width $n+1$ and no internal biases.
\end{proof}

The theorem rules out a polynomial-cost flattening into a single hidden
layer, even when arbitrarily many hidden units and arbitrarily precise
real coefficients are allowed. It does not rule out a polynomial-size
constant-depth threshold representation; the displayed target already
has one. Nor does this unit-margin exact-representation lower bound by
itself lower-bound the error of approximate sign representations.

In fact Corollary~\ref{cor:fixed-intrinsic-depth} applied to the explicit
depth-three construction gives the original unrestricted posterior a
sample bound polynomial in $n,L,\varepsilon^{-1},\log(2/\delta)$ for
every $L\ge3$, every actual width $N\ge C_L(\mathrm{IP}_n)$, and every
$0<\sigma^2\le N^{-(1+\zeta)}$, for each fixed $\zeta>0$.
An exponential shallow norm is therefore compatible with efficient
statistical learning by this deep posterior. It cannot on its own
supply the missing unrestricted-depth counterexample.

\section{Excess Gaussian energy and intrinsic product complexity}
\label{sec:intrinsic-product-energy}

Define a function-level complexity, allowing all finite architectures,
by
\[
 \mathfrak P(f)=\inf\left\{
   \prod_j\|W^j\|_F:
   f(x)z_w(x)\ge1\ \hbox{on the entire cube}
   \right\}.
\]
The infimum permits arbitrary finite depth and width, retains tanh
hidden layers and a linear readout, and allows biases only through
the augmented input. It is finite, for example by a truth-table
shallow construction. It is also at least one: the network extends
oddly to the full augmented cube $\{-1,1\}^{d+1}$ and has absolute
margin at least one there. This nonconstant odd Boolean extension
has a bichromatic cube edge. Along that edge the product of layer
operator norms, and hence the Frobenius product, must be at least one.

Let $J_{L,N}(f)$ be the infimum of the Gaussian energy $\kappa_N$
over width-$N$, depth-$L$ unit-margin representations of $f$, with
$J_{L,N}(f)=+\infty$ if none exists. This is a deterministic energy
quantity, not a finite-temperature posterior probability.

\begin{theorem}[Intrinsic product complexity characterizes excess energy]
\label{thm:intrinsic-product-energy}
For every $d\ge1$, $L\ge2$ and $N\ge1$,
\[
 J_{L,N}(f)\ge L+\log d+2\log\mathfrak P(f).
\]
Conversely, there are universal polynomials $p,q$ and a universal
constant $c$ such that
\[
 L\ge p(\mathfrak P(f))+c\log(d+1),\qquad
 N\ge q(d,\mathfrak P(f))
\]
imply
\[
 J_{L,N}(f)\le L+c\log(2d\mathfrak P(f)).
\]
The upper-bound witness can be chosen with all physical weights at
most $q(d,\mathfrak P(f))$. The polynomials do not depend on any
supplied representation's depth, width, or coefficient precision.
\end{theorem}

\begin{proof}
The product of the $L$ weighted layer energies is
$d\prod_j\|W^j\|_F^2$. Arithmetic--geometric mean and
$e^u\ge1+u$ therefore give
\[
 \kappa_N\ge L\left(d\prod_j\|W^j\|_F^2\right)^{1/L}
 \ge L+\log d+2\log\mathfrak P(f).
\]
This proves the lower bound without any width or depth restriction.

Choose a finite representation of product norm less than
$2\mathfrak P(f)$. The depth-independent normal form supplies a
unit-margin representation whose depth $D_0$, width $C$, physical
weights and product norm $P_0$ are all polynomial in $\mathfrak P(f)$.
Apply the product-target construction of
Theorem~\ref{thm:product-target-log-padding}. Its width threshold $M$
and transmission parameter $T$ are polynomial in $d,C,D_0,P_0$,
hence polynomial in $d,\mathfrak P(f)$, and do not depend on the
subsequently chosen learner depth. For
$L\ge D_0+\lceil\log T\rceil$, that construction has energy at most
\[
 L+15+32\log T.
\]
If the small-prefix linearization in that proof removes every hidden
activation, its separate halfspace learning conclusion is not needed
here. The collapsed affine map has norm at most $2P_0$ and margin
greater than one. Apply tanh to that map and append readout $2$.
This is a width-one shallow unit-margin representation with weights
polynomial in $P_0$. The shallow scalar-transmission construction
therefore supplies, at every sufficiently large polynomial width and
every $L\ge2+O(\log(d+P_0))$, energy
$L+O(\log(2dP_0))$. It also has polynomially bounded physical weights.
Thus the affine branch has the same claimed upper-bound form.

Its physical weights are bounded by a polynomial width threshold.
Since $\log T=O(\log(d+1)+\log\mathfrak P(f))$, increasing universal
polynomials and constants gives the stated conditions and upper
bound. No attainment of the infimum defining $\mathfrak P(f)$ was
assumed.
\end{proof}

\begin{corollary}[Learning from intrinsic product complexity]
\label{cor:intrinsic-product-learning}
For each fixed $\zeta>0$, the unrestricted Gaussian posterior has a
sample bound polynomial in
\[
 d,\mathfrak P(f),L,\varepsilon^{-1},\log(2/\delta)
\]
at every $L\ge p(\mathfrak P(f))+c\log(d+1)$, simultaneously for
all actual widths $N\ge C_L(f)$ and all
$0<\sigma^2\le N^{-(1+\zeta)}$.
\end{corollary}

\begin{proof}
Use the same depth-independent normal form and apply the all-width
product-target learning theorem. Its intrinsic parameters are all
polynomial in $\mathfrak P(f)$. Its finite-width branch covers the
interval between the actual minimum width and the constructed
polynomial-width witness threshold.
\end{proof}

Thus, once depth and width are sufficiently large, the excess above
the depth baseline is, up to universal constants, the logarithm of
intrinsic product complexity together with the input normalization
cost $\log d$. This does not identify $\mathfrak P(f)$ with ordinary
circuit size. The normal form proves that ordinary Boolean description
size is polynomial in $\mathfrak P(f)$; a converse polynomial bound
for arbitrary small deep circuits remains the central unresolved
representation question. The learning corollary also retains its
displayed learner-depth threshold.

\section{Depth alignment and majority combinations}
\label{sec:pf-majority-closure}

A possible route from distribution-dependent approximation to exact
representation uses minimax and a majority of approximants. The unresolved
step is combining their different nonlinear computations at polynomial
product cost. Unequal depths alone are not an obstruction: a scalar
readout can be padded at constant product cost. A stronger conclusion
holds when the approximants already share their nonlinear feature map.

\begin{lemma}[Constant-product scalar padding]
\label{lem:pf-scalar-padding}
A depth-$D$ tanh network with global signed margin at least one and
Frobenius product $P$ has, for every integer $T\ge1$, a depth-$(D+T)$
realization of the same Boolean function, with global signed margin
at least one and product at most $3P/2$. No new biases are required.
\end{lemma}

\begin{proof}
Write $z$ for the old readout and $R=\sqrt{d+1}$.
Zero preservation and one-Lipschitzness give $|z(x)|\le PR$.
Put $M=\max\{1,PR\}$ and $\eta=(M\sqrt T)^{-1}$, and replace the
readout by
\[
 h_1=\tanh(\eta z),\qquad h_{s+1}=\tanh(h_s),\qquad
 z_{\rm new}=\frac{3h_T}{2\eta}.
\]
The hidden absolute values are at most $T^{-1/2}$. The inequality
$\tanh u\ge u-u^3/3$ for $u\ge0$ and Bernoulli's inequality imply
\[
 |h_T|\ge\eta|z|\left(1-\frac1{3T}\right)^T
           \ge\frac23\eta|z|.
\]
Oddness preserves the sign. The original readout norm is multiplied
by $\eta$, the next $T-1$ scalar weights equal one, and the new
readout weight is $3/(2\eta)$. Their product changes by only $3/2$.
\end{proof}

\begin{proposition}[A disjoint-branch product bound]
\label{prop:pf-parallel-holder}
Suppose $m$ independent branches of common weighted depth $D$ are
stacked at their first layer, remain disjoint internally, and are
combined linearly at the readout. Let $a_{i\ell}>0$ be branch $i$'s
layer norms, including its voting coefficient in its final norm, and
put $P_i=\prod_\ell a_{i\ell}$. The complete product satisfies
\[
 P_{\rm total}\ge
       \left(\sum_{i=1}^m P_i^{2/D}\right)^{D/2}.
\]
In particular, averaging $m$ nonconstant globally unit-margin branches
has product at least $m^{D/2-1}$, regardless of their internal widths
or neuron-cloning choices.
\end{proposition}

\begin{proof}
Each full matrix has squared Frobenius norm $\sum_i a_{i\ell}^2$.
H\"older's inequality gives
\[
 \sum_i\prod_\ell a_{i\ell}^{2/D}
 \le\prod_\ell\left(\sum_i a_{i\ell}^2\right)^{1/D}.
\]
Raise this to power $D/2$. Every nonconstant unit-margin Boolean
branch has product at least one, by a bichromatic cube edge.
Including its averaging coefficient gives $P_i\ge1/m$.
These are the products of the actual branches after any cloning,
which proves the last assertion.
\end{proof}

Uniform cloning of a whole layer preserves its branch product, but
neuron-by-neuron cloning can reduce it. The proposition uses the
resulting products and their intrinsic margin lower bounds, so it does
not require the stronger, false invariance claim. It is an obstruction
to a disjoint architecture, not a lower bound on every implementation
of the resulting majority function.

\begin{theorem}[Voting with a shared nonlinear feature map]
\label{thm:pf-shared-feature-voting}
Let $h(x)$ be one common tanh feature map with prefix product $S$.
Suppose $z_i(x)=v_i^Th(x)$ satisfy $|z_i(x)|\ge1$ everywhere,
let $g_i=\operatorname{sign}z_i$, and put $P_i=S\|v_i\|_2$.
If $\lambda_i\ge0$, $\sum_i\lambda_i=1$, and
\[
 f(x)\sum_i\lambda_i g_i(x)\ge\gamma>0
 \quad\hbox{at every input},
\]
then
\[
 \mathfrak P(f)\le
 \frac{\log(4/\gamma)}{\gamma}\sum_i\lambda_iP_i.
\]
A concrete finite-width realization attains at most twice the displayed
bound and adds just one weighted layer relative to the networks $z_i$.
It uses no new biases.
\end{theorem}

\begin{proof}
Set $A=\tfrac12\log(4/\gamma)$. Since $|z_i|\ge1$,
$|\tanh(Az_i)-g_i|\le\gamma/2$. Therefore
\[
 z_{\rm new}=\frac2\gamma\sum_i\lambda_i\tanh(Av_i^Th)
\]
has global signed margin at least one. Clone its $i$th new hidden
neuron $c_i$ times and divide its outgoing coefficient by $c_i$;
this leaves the function exactly unchanged. Its product is
\[
 \frac{2AS}{\gamma}
 \sqrt{\left(\sum_i c_i\|v_i\|_2^2\right)
       \left(\sum_i\lambda_i^2/c_i\right)}.
\]
Omit zero coefficients and take
$c_i=\lceil t\lambda_i/\|v_i\|_2\rceil$.
As $t\to\infty$, the square root tends to
$\sum_i\lambda_i\|v_i\|_2$. This gives the intrinsic infimum bound;
a sufficiently large finite $t$ gives its factor-two version.
The vectors $v_i$ are nonzero by their margin assumption.
\end{proof}

For an odd unweighted majority, $\gamma=1/m$, so the bound is
$O(m\log(m+1)\max_iP_i)$. For a fixed positive voting margin it is
$O(\max_iP_i)$, independently of the number of voters.

\begin{corollary}[Weak-to-exact representation with shared features]
\label{cor:pf-shared-feature-minimax}
Fix the feature map $h$ and a budget $P$. Suppose that for every
probability distribution $\mu$ on the cube there is a readout $v_\mu$
with global margin $|v_\mu^Th|\ge1$, product
$S\|v_\mu\|_2\le P$, and correlation
\[
 \mathbb E_\mu\bigl[f(x)\operatorname{sign}(v_\mu^Th(x))\bigr]
 \ge\gamma>0.
\]
Then
\[
 \mathfrak P(f)\le P\frac{\log(4/\gamma)}{\gamma}.
\]
Consequently the depth-independent normal form gives a finite
polynomial-size representation in this product bound.
\end{corollary}

\begin{proof}
Retain one admissible readout for every distinct Boolean sign pattern
in the specified class. The class is finite because the input set is
finite. The finite minimax theorem gives a probability distribution on
these readouts whose signed vote is at least $\gamma$ at every input.
Apply Theorem~\ref{thm:pf-shared-feature-voting}.
\end{proof}

The shared feature map is a substantive hypothesis. It is not supplied
by allowing a different arbitrary low-product network for each input
distribution. General polynomial closure under majority remains open.
Preserving the original shallow depth is already false: the functions
$q_j=\operatorname{sign}(2\sum_i u_iv_i-2j+1)$ in
Section~\ref{sec:shallow-product-separation} have shallow polynomial
product on shared pairwise-AND features, while the alternating votes
$(-1)^jq_j$ and the optional constant vote form an odd majority equal
to $\mathrm{IP}_n$. The latter has exponential shallow product, though
one additional layer gives polynomial product.

\section{Depth-independent input editing}
\label{sec:primary-input-editing}

A shallow computation placed before an arbitrarily deep network can
modify its Boolean function without carrying extra features through
the old depth. We use this observation to change finitely many labels,
and to combine a network with a predicate of small Fourier variation.
The product bounds depend nearly linearly on the old product, with
no dependence on the old depth.

Write $\bar x=(1,x)$ and let $z(u)$ denote the extension of a network
to arbitrary augmented inputs $u\in\mathbb R^{d+1}$. If its Frobenius
product is $P$, then
\[
 z(-u)=-z(u),\qquad |z(u)-z(v)|\le P\|u-v\|_2.
\]
The first identity uses tanh hidden activations, a linear readout, and
the absence of biases after the augmented input.

\begin{lemma}[Vector-output cloning and precomposition]
\label{lem:primary-vector-cloning}
For a shallow vector map
\[
 T(u)=\sum_j b_j\tanh(a_j^{\mathsf T}u),
 \qquad V=\sum_j\|a_j\|_2\|b_j\|_2,
\]
finite neuron cloning gives an identical map whose two matrix
Frobenius norms have product at most $2V$. If this map is placed
before a network of product $P$, its linear output can be merged
with the old first matrix. The resulting product is at most $2PV$,
and only one hidden layer is added.
\end{lemma}

\begin{proof}
Cloning neuron $j$ exactly $c_j$ times and dividing its outgoing
vector by $c_j$ gives product
\[
 \left\{\left(\sum_jc_j\|a_j\|_2^2\right)
              \left(\sum_j\|b_j\|_2^2/c_j\right)\right\}^{1/2}.
\]
After omitting zero terms, choose
$c_j=\lceil t\|b_j\|_2/\|a_j\|_2\rceil$ and let $t$ increase.
The expression tends to $V$, so a finite choice is within a factor
two. If the map's matrices are $A,B$ and the old first matrix is
$W_1$, use $\|W_1B\|_F\le\|W_1\|_F\|B\|_F$.
\end{proof}

\begin{theorem}[Changing finitely many Boolean labels]
\label{thm:primary-finite-label-editing}
Suppose $g(x)z(\bar x)\ge1$ everywhere on the $d$-dimensional
Boolean cube, $d\ge1$, and the supplied representation has product
$P\ge1$. Let $E$ contain $k\ge1$ distinct inputs, and set
$f=-g$ on $E$ and $f=g$ elsewhere. There is a unit-margin
representation of $f$ with one additional hidden layer and product
at most
\[
 CP\left\{d+1+kd\sqrt{d+1}\,
                   \log\bigl(8Pk\sqrt{d+1}\bigr)\right\}.
\]
Here $C$ is universal. The new prefix has finite width polynomial
in $d,k,\log(P+2)$. The bound is independent of the old depth and
width. The same inequality holds with $P=\mathfrak P(g)$, after
enlarging $C$.
\end{theorem}

\begin{proof}
Set $R=\sqrt{d+1}$, $h=1-2\mathbf1_E$, and
\[
 T(x)=h(x)\bar x
     =\bar x-2\sum_{e\in E}\bar e\,\mathbf1_{\{x=e\}}.
\]
Oddness gives $z(T(x))=h(x)z(\bar x)$, so its signed margin for
$f$ is at least one, including when $g$ is constant.

For every $e\in E$, the affine form
$s_e(x)=e^{\mathsf T}x-d+1$ is $+1$ at $e$ and at most $-1$
elsewhere. Put
\[
 \delta=\frac1{4PkR},\qquad
 \lambda=\tfrac12\log(8PkR),\qquad
 t_e(x)=\tanh(\lambda s_e(x)).
\]
Then $|t_e-(2\mathbf1_{\{x=e\}}-1)|\le\delta$. Consequently
\[
 \widetilde T(x)=\bar x-\sum_{e\in E}\bar e\,(1+t_e(x)),
 \qquad
 \|\widetilde T(x)-T(x)\|_2\le kR\delta=\frac1{4P}.
\]
It follows that $2z(\widetilde T(x))$ has signed margin at least
$3/2$ for $f$.

This vector map has a shallow tanh realization. Write
$\alpha=\operatorname{atanh}(1/2)$, $\tau=\tanh(1)$, and let
$e_0,\ldots,e_d$ denote the standard coordinate vectors. Use one
constant neuron with incoming row $\alpha e_0^{\mathsf T}$ and
outgoing vector $2(e_0-\sum_{e\in E}\bar e)$; $d$ identity neurons
with incoming rows $e_i^{\mathsf T}$ and outgoing vectors
$e_i/\tau$; and the $k$ point neurons with incoming rows
$\lambda(1-d,e^{\mathsf T})$ and outgoing vectors $-\bar e$.
The identity coordinates are exact on the cube, since
$\tanh(x_i)=\tau x_i$. The variation of this vector map is at most
\[
 V\le2\alpha(1+kR)+d/\tau+k\lambda dR,
\]
using $\sqrt{(d-1)^2+d}\le d$. Apply
Lemma~\ref{lem:primary-vector-cloning}, merge the prefix, and double
the old readout. This gives product at most $4PV$, as required.

For the width assertion, the explicit choice
$t=\max_j\|a_j\|_2/\|b_j\|_2$ in the cloning proof gives product
at most $\sqrt2V$: the ceiling contributes at most
$t\|a_j\|_2\|b_j\|_2$ to the incoming squared norm. For the
displayed nonzero rows and columns, both this $t$ and every ratio
$\|b_j\|_2/\|a_j\|_2$ are polynomial in $d,k,\lambda$.
Thus the number of clones is polynomial as claimed.

For intrinsic complexity, take a finite representation of product
at most $2\mathfrak P(g)$ and increase the universal constant.
No attainment of the defining infimum is needed.
\end{proof}

\begin{lemma}[Fourier input editing]
\label{lem:primary-fourier-input-editing}
Suppose $T:\{-1,1\}^d\to\mathbb R^r$ has vector Fourier
expansion $T=\sum_S b_S\chi_S$ and variation
$A=\sum_S\|b_S\|_2>0$. For $0<\eta<1$, a shallow tanh vector
map approximates $T$ uniformly in Euclidean norm to error $\eta$
with product at most
\[
 C(d+1)^2A\log\left(\frac{2(d+1)\max\{A,1\}}{\eta}\right).
\]
If a network of product $P\ge1$ has signed margin one after exact
composition with $T$, the composed Boolean function has a
representation with one additional hidden layer and product at most
\[
 CP(d+1)^2A\log\bigl(8(d+1)\max\{A,1\}P\bigr).
\]
\end{lemma}

\begin{proof}
For $S\ne\varnothing$, write $s=|S|$ and let $t$ be the number
of negative selected input coordinates. The exact identity
\[
 \chi_S(x)=b_0+\sum_{j=1}^{s}(-1)^j
                    \operatorname{sign}(2t-2j+1),
 \qquad b_0=1+\sum_{j=1}^{s}(-1)^j\in\{0,1\},
\]
represents parity using affine thresholds of integer margin one.
Replacing their signs by tanh at steepness
$\tfrac12\log(2s/\delta)$ approximates parity to error $\delta$.
Its scalar variation is at most
$C(d+1)^2\log(2(d+1)/\delta)$. A constant neuron supplies $b_0$
exactly, and also supplies the empty-set parity.

Take $\delta=\min\{1/2,\eta/(2\max\{A,\eta\})\}$ and
multiply the outgoing coefficients of the $S$th parity approximation
by $b_S$. The vector error is at most $A\delta\le\eta$, and
the sum of vector variations has the asserted bound.
Lemma~\ref{lem:primary-vector-cloning} supplies the matrix-product
bound. For composition take $\eta=1/(4P)$, use the old network's
$P$-Lipschitz property, merge the prefix, and double the final
readout.
\end{proof}

\begin{corollary}[Multiplication by a low-variation predicate]
\label{cor:primary-predicate-modulation}
Let $h$ be Boolean with Fourier variation
$H=\sum_S|\widehat h(S)|$. Then
\[
 \mathfrak P(gh)
 \le C\mathfrak P(g)(d+1)^3H
       \log\bigl(8(d+1)^2H\mathfrak P(g)\bigr).
\]
In particular, multiplying by any signed parity has a cost
polynomial in $d$ and nearly linear in $\mathfrak P(g)$,
independently of the depth of a supplied representation of $g$.
The same order of bound holds for $g\mathbin{\mathrm{AND}}h$
and $g\mathbin{\mathrm{OR}}h$, with $+1$ representing true.
\end{corollary}

\begin{proof}
For multiplication use $T(x)=h(x)\bar x$. Each coordinate has
Fourier variation $H$, so its vector Fourier variation is at most
$(d+1)H$. Oddness of the old network and
Lemma~\ref{lem:primary-fourier-input-editing} prove the bound.
A signed parity has $H=1$.

For AND, if $g(a)=-1$ at some cube point $a$, use
\[
 T(x)=\left(1,\frac{1+h(x)}2x+\frac{1-h(x)}2a\right).
\]
It equals $\bar x$ when $h=+1$ and $\bar a$ when $h=-1$,
and has vector Fourier variation at most $C(d+1)(H+1)$.
If $g$ is the constant $+1$, the answer is $h$, which has a direct
shallow representation by the parity construction in the preceding
proof. The constant $-1$ case is immediate. OR uses a positive
witness and the reversed selector. These are existence statements;
no efficient procedure for finding the witness is required.
\end{proof}

These results do not prove closure under majority of arbitrary
polynomial-product networks. A general deep predicate need not
have small Fourier variation. Constructing the vector map
$h(x)(1,x)$ from such a predicate with polynomial product remains
an identity-augmentation problem: carrying $x$ through all its
layers by parallel coordinates incurs the existing depth cost.
Likewise, finite-label editing supplies an exact interpolation
witness but does not compare it with inaccurate interpolants, so
it is not by itself a noisy-label generalization theorem.


\section*{Analytic and representation foundations}
The following sections provide the full supporting arguments. They are
included to make this paper independent of earlier drafts in this
research program.

\section{Unrestricted learning at bounded widths}
\label{sec:finitewidth-unrestricted}

A robustness restriction is unnecessary when the learner's width is itself
polynomially bounded. The following elementary argument retains the original
all-input clean margin and Gaussian conditioning conventions, and is uniform
in the Gaussian scale. It also covers analog target networks at their original
depth, without digital simulation.

Let $P_N=N(d+2)+(L-2)N^2$ be the parameter count, and let $\mu_{N,t}$ denote
the raw Gaussian with variances $t^2/d,t^2/N,t^2/N^2$, respectively. Thus
$\kappa_N=(t^2/N)\chi^2_{P_N}$ under $\mu_{N,t}$. Fix integers
$1\le C\le M$ and suppose, for some $D\ge1$, that
\begin{equation}\label{eq:finitewidth-targetmass}
 \mu_{N,t}\{w:f(x)z_w(x)>1/2\text{ for every }x\}
 \ge \exp(-DN/t^2).
\end{equation}
This holds for $C\le N\le M$ and $0<t\le1$. Only this target-mass
hypothesis is needed; posterior networks may be fragile.

\begin{theorem}\label{thm:finitewidth-unrestricted}
For $0<\varepsilon,\delta<1$, set
\begin{align*}
 B&=4D+2[d+2+(L-2)M]\log2+4\log(2/\varepsilon),\\
 A&=\sqrt{MB},\qquad T=d+1+M(L-1),\qquad
 q=[4T(MA)^{L-1}]^{-1}.
\end{align*}
If
\[
 n\ge\frac2\varepsilon\left[
 \log M+P_M\log(5A/q)+\log(1/\delta)\right],
\]
then, with probability at least $1-\delta$ over independent examples from any
input distribution, the unrestricted Gaussian posterior has Gibbs error at
most $\varepsilon$, simultaneously for every $C\le N\le M$ and $0<t\le1$.
Here ``unrestricted'' retains the original conditioning on
$|z_w(x)|>1/2$ for all inputs, but imposes no internal robustness condition.
The same sample event works after any number of additional correctly labeled
examples. In particular, if $D$ and $M$ are polynomial in $d,L,C$, the sample
bound is polynomial in $d,L,C,\varepsilon^{-1},\log(1/\delta)$, uniformly for
all $\gamma\ge1$ with $t^2=N^{1-\gamma}$.
\end{theorem}

\begin{proof}
For every sample, the event used to condition the raw Gaussian contains the
target event in~\eqref{eq:finitewidth-targetmass}. The chi-square exponential
moment therefore gives
\[
 Q_{S,N,t}\{\kappa_N>B\}
 \le\exp\left(\frac{DN}{t^2}+\frac{P_N}2\log2-
                       \frac{NB}{4t^2}\right)
 \le\varepsilon/2.
\]
The last inequality uses $P_N/N\le d+2+(L-2)M$, $N\ge1$, and $t\le1$.
This bound holds for every sample, without a probabilistic qualification.

If $\kappa_N\le B$, every physical parameter has magnitude at most $A$.
Round each to the nearest multiple of $q$, using at most $5A/q$ grid values
per coordinate. For the maximum coordinatewise hidden-state discrepancies,
boundedness and the Lipschitz property of tanh give
\[
 e_1\le(d+1)q,\qquad e_s\le NAe_{s-1}+Nq,\qquad
 |z-z_{\rm grid}|\le NAe_{L-1}+Nq.
\]
Consequently, uniformly over Boolean inputs,
\[
 |z-z_{\rm grid}|\le q[d+1+M(L-1)](MA)^{L-1}=1/4.
\]
Rounding thus preserves the entire Boolean function of every posterior
network with energy at most $B$.

The union of these grid classifiers over $C\le N\le M$ has logarithmic
cardinality at most $\log M+P_M\log(5A/q)$. A consistency union bound gives
one sample event on which every consistent grid classifier has true error
at most $\varepsilon/2$. Every good-energy posterior draw has the same
function as such a classifier; the remaining posterior mass is at most
$\varepsilon/2$. This proves the Gibbs bound. The finite class is independent
of $t$, and consistency only becomes more restrictive as examples are added.
\end{proof}

The logarithm of the grid size satisfies
\[
 \log(5A/q)=\log(20AT)+(L-1)\log(MA).
\]
The exponentially small parameter precision costs only a polynomial in
$L$ in this logarithm. The unresolved unrestricted problem is therefore
uniformity at learner widths \emph{larger than every polynomial in the
specified target complexity and depth}; this theorem does not show those
larger widths actually require more samples.

\begin{corollary}[Same-depth bounded-weight analog targets]
Suppose a width-$C$, depth-$L$ tanh network realizes $f$ with
$f(x)z(x)\ge1$ on every Boolean input and weights of magnitude at most
$R\ge1$. Fix $M\ge C$ and set
\begin{align*}
 K&=2[d(d+1)+(L-1)C^2]R^2,\\
 a&=[4(d+1+M(L-1))(MR)^{L-1}]^{-1},\\
 D&=\max\{1,K+[d+2+(L-2)M][2+\log(1/a)]\}.
\end{align*}
Then~\eqref{eq:finitewidth-targetmass} holds with this $D$.
Thus Theorem~\ref{thm:finitewidth-unrestricted} supplies a polynomial sample
bound at the original depth for every $C\le N\le M$, with dependence on
$R$ only through $\log R$ in the final sample bound. No robustness condition
on the target implementation or prior support is needed. A target with the
original strict margin $>1/2$ satisfies the displayed hypothesis after
doubling its readout.
\end{corollary}

\begin{proof}
Clone the $C$ neurons into balanced nonempty groups of sizes $n_i$ summing
to $N$. Copy first-layer rows and replace internal and readout weights by
$W_{ij}/n_j$ and $v_i/n_i$. This preserves the exact function. Since
$\max_i n_i/\min_i n_i\le2$, its energy is at most $K$, and its weights
remain bounded by $R$. A physical-coordinate perturbation of magnitude
at most $a$ changes every output by at most $1/4$, by the recurrence above.

Use isotropic coordinates
$y=(\sqrt d\,W^1,\sqrt N\,W^{\rm internal},Nv)$, so
$y\sim\mathcal N(0,t^2I)$ and $\|y_*\|^2\le NK$ for the cloned center.
The standard-normal-coordinate box
$|y_j/t-y_{*,j}/t|\le a$ lies within the physical box, because all coordinate
scalings are at least one and $t\le1$. Its probability is at least
\begin{align*}
 &\left(\frac{2a}{\sqrt{2\pi}}\right)^{P_N}
 \exp\left[-\frac12\sum_j(|y_{*,j}|/t+a)^2\right]\\
 &\qquad\ge\exp[-NK/t^2-P_N(2+\log(1/a))]
 \ge\exp(-DN/t^2).
\end{align*}
Here $(u+a)^2\le2u^2+2a^2$ and $a\le1$ suffice. Every network in this box
has target margin at least $3/4$.
\end{proof}

This is an elementary finite-parameter learning result, rather than a new
claim about the general VC dimension of neural networks. Its role is to
isolate which width regime remains open under the exact prior conventions.

\section{A width-independent tanh bound without robustness}
\label{sec:productnorm}

We first establish the width-independent function-class bound used to
control the low-energy signal extracted from a posterior draw. The proof
accounts for the fact that tanh is not positively homogeneous.

The normalized-activation and exponential-moment method is closely related
to Nakhleh--Nowak's treatment of deep neural variation spaces
\cite{nakhleh2026}, which explicitly uses $\tanh(st)/s$ and slope
discretization, and to the exponential peeling argument of
Golowich--Rakhlin--Shamir~\cite{golowich2018}. We give a self-contained
product-Frobenius version for the present Gaussian-energy conventions;
we do not claim the underlying proof technique is new.

\begin{theorem}[Product-Frobenius bound for tanh]\label{thm:tanh-product}
Let $\mathcal F_{L,P}$ contain all scalar-output tanh networks with $L-1$
hidden layers, arbitrary finite widths, a linear readout, and product of
layer Frobenius norms at most $P$. Biases are absent except through an
augmented input. On every sample $x_1,\ldots,x_n$ with $\|x_i\|\le R$,
\[
 \widehat{\mathfrak R}_n(\mathcal F_{L,P})
 \le \frac{PR}{\sqrt n}
 \left[2+\sqrt{2(L-1)\log\bigl(3+n(L-1)^2\bigr)}\right].
\]
The bound is independent of every hidden width and of the individual layer
norms apart from their product.
\end{theorem}

\begin{proof}
If $R=0$, every sampled input and every network output on the sample is
zero, so the bound is immediate. Assume $R>0$ below.
Define $\phi_s(u)=\tanh(su)/s$ for $s>0$, $\phi_0(u)=u$, and
$\phi_\infty(u)=0$. These functions are odd, $1$-Lipschitz, and zero at zero.
For finite $s,t\ge0$ and $r\ge0$,
\[
 \phi_s(ru)=r\phi_{sr}(u),\qquad
 |\phi_s(u)-\phi_t(u)|\le u^2|s-t|,
 \qquad |\phi_s(u)|\le1/s\quad(s>0).
\]
The derivative bound follows by differentiating
$\phi_s(u)=u\int_0^1\operatorname{sech}^2(suv)\,dv$.

Let $\mathcal G_0$ contain the identity feature map. Inductively allow
arbitrary finite-dimensional features
\[
 g_j(x)=\phi_{s_j}(a_j^T h(x)),\quad
 h\in\mathcal G_{\ell-1},\quad \sum_j\|a_j\|^2\le1,
 \quad s_j\in[0,\infty].
\]
All these feature vectors have norm at most $R$ on the sample. Put
$B_\ell(\epsilon)=\sup_{g\in\mathcal G_\ell}
 \|\sum_i\epsilon_i g(x_i)\|$ for independent Rademacher signs.
Writing $a_j=r_j u_j$ and using the scaling identity shows
\[
 B_\ell\le
 \sup_{h,\,\|u\|\le1,\,s\in[0,\infty]}
 \left|\sum_i\epsilon_i\phi_s(u^T h(x_i))\right|.
\]
For any $\eta>0$, uniformly on $[-R,R]$ the slope family admits an
$\eta/n$ approximation using at most
$J=3+n^2R^2/\eta^2$ choices: truncate $s>n/\eta$ to the zero function,
and round slopes below that threshold down at mesh $\eta/(nR^2)$.
The error in the Rademacher sum is at most $\eta$.

We use the convex contraction principle
\[
 \mathbb E G\!\left(\sup_{t\in T}\sum_i\epsilon_i\psi_i(t_i)\right)
 \le
 \mathbb E G\!\left(\sup_{t\in T}\sum_i\epsilon_i t_i\right)
\]
for convex nondecreasing $G$ and $1$-Lipschitz $\psi_i$ vanishing at zero.
For completeness, the single-coordinate proof compares
$X_\pm=\sup_t(a_t\pm\psi(t_k))$ and
$Y_\pm=\sup_t(a_t\pm t_k)$. One has
$\max(Y_+,Y_-)\ge\max(X_+,X_-)$ by $|\psi(t)|\le|t|$.
An approximately maximizing pair for $X_+,X_-$ also gives
$Y_++Y_-\ge X_++X_-$, by the Lipschitz property.
Convex monotonicity therefore gives
$G(Y_+)+G(Y_-)\ge G(X_+)+G(X_-)$; iterate over coordinates.
Bounded separable classes follow by approximation from finite classes.
In our application, $u$ ranges symmetrically and $\phi_s$ is odd, so the
absolute supremum is the same as the one-sided supremum.

Apply this principle with $G(t)=e^{\lambda t}$ to each grid slope, then
bound the exponential of a maximum by the sum of the exponentials:
\[
 \mathbb E e^{\lambda B_\ell}
 \le e^{\lambda\eta}J\,\mathbb E e^{\lambda B_{\ell-1}}.
\]
Furthermore, $\mathbb E B_0\le R\sqrt n$, and bounded differences gives
$\mathbb E e^{\lambda B_0}
 \le e^{\lambda R\sqrt n+\lambda^2nR^2/2}$.
Writing $H=L-1$, iteration and Jensen yield
\[
 \mathbb E B_H\le H\eta+R\sqrt n+
             \frac{H\log J}{\lambda}+\frac{\lambda nR^2}{2}.
\]
Take $\eta=R\sqrt n/H$ and optimize $\lambda$ to obtain
$\mathbb E B_H\le R\sqrt n[2+\sqrt{2H\log(3+nH^2)}]$.

Finally, for a particular original network, write
$M_\ell=\|W^\ell\|_F$ and $S_\ell=\prod_{j\le\ell}M_j$.
Dividing its $\ell$th hidden vector by $S_\ell$ changes its recursion to
$\phi_{S_\ell}((W^\ell/M_\ell)g^{\ell-1})$.
Thus its normalized hidden features belong to the enlarged classes above,
and its scalar output is at most $P$ times a norm-one linear readout of
$\mathcal G_H$. Zero-norm layers give the zero function and require no
normalization. This proves the theorem. All suprema are measurable by
countable approximation of finite-dimensional parameter spaces and then
union over widths.
\end{proof}

For the Gaussian energy in our model, the product of the $L$ weighted
squared layer norms is $dP^2$. Consequently
\[
 P\le d^{-1/2}(\kappa_N/L)^{L/2}
 \le d^{-1/2}e^{B/2}\qquad\text{if }\kappa_N\le L+B.
\]
Applying the theorem to the ramp loss
$\ell(yz)=\min\{1,\max\{0,1-2yz\}\}$ gives the following explicit bound.

\begin{corollary}[All widths at bounded excess energy]\label{cor:excess-risk}
Fix $B\ge0$. With probability at least $1-\delta$, every sample-consistent
depth-$L$ network, at any finite width, with training margins at least $1/2$
and energy $\kappa_N\le L+B$ has sign error at most
\[
 \frac{4\sqrt2 e^{B/2}}{\sqrt n}
 \left[2+\sqrt{2(L-1)\log\bigl(3+n(L-1)^2\bigr)}\right]
 +\sqrt{\frac{\log(1/\delta)}{2n}}.
\]
No internal robustness assumption is imposed.
\end{corollary}

\begin{proof}
Augmented Boolean inputs have norm $\sqrt{d+1}$, so $PR\le\sqrt2e^{B/2}$.
The ramp loss is bounded by one and is $2$-Lipschitz. Symmetrization,
scalar contraction, and bounded differences give population loss at most
empirical loss plus four times the network Rademacher bound and the
stated confidence term. Empirical loss is zero and population loss
upper-bounds sign error.
\end{proof}

\begin{lemma}[A scalar transmission construction]\label{lem:scalar-transmission}
For every real $D>0$ and integer $m\ge2$, there are positive scalar weights
$\alpha,b_2,\ldots,b_{m-1},c$ such that
\[
 r_1=\tanh\alpha,\quad r_s=\tanh(b_s r_{s-1}),\quad
 cr_{m-1}=1/2,
\]
and, writing $D'=\max\{D,4\}$,
\[
 D\alpha^2+\sum_{s=2}^{m-1}b_s^2+c^2
 \le(m+12)(D'/4)^{1/m},\qquad c^2=(D'/4)^{1/m}.
\]
\end{lemma}

\begin{proof}
Set $q=(D'/4)^{1/(2m)}$, $k_s=\min\{s,m-s\}$, and
$r_s=q^s/(\sqrt{D'}k_s)$. Then $0<r_s\le1/(2qk_s)\le1/2$ and
$r_{m-1}=1/(2q)$. Choose
$\alpha=\operatorname{atanh}r_1$,
$b_s=\operatorname{atanh}(r_s)/r_{s-1}$ and $c=q$.
For $0<u\le1/2$ the power series gives
$(\operatorname{atanh}(u)/u)^2\le1+u^2$.
Writing $R_s=k_{s-1}/k_s$, the reciprocal pairs
$k/(k+1),(k+1)/k$ satisfy
\[
 \left(\frac{k+1}{k}\right)^2+
 \left(\frac{k}{k+1}\right)^2-2\le4/k^2.
\]
Consequently $\sum_{s=2}^{m-1}R_s^2\le m-2+7$,
$R_s\le2$, and $\sum_{s=1}^{m-1}k_s^{-2}<4$.
Using $q^{2m}=D'/4$ therefore gives
\[
 D\alpha^2\le q^2+1/4,\qquad
 \sum_{s=2}^{m-1}b_s^2\le q^2(m-2+7)+4,\qquad c^2=q^2.
\]
Their sum is at most $(m+12)q^2$. Empty internal sums at $m=2$
satisfy the same estimates.
\end{proof}


\section{Unrestricted finite-variance learning for \texorpdfstring{$\gamma>1$}{gamma greater than one}}
\label{sec:gaussian-distance}

There is a stronger bridge than concentration of the full Frobenius energy.
A typical Gaussian bulk has small operator norms when $\gamma>1$, even when
its total Frobenius energy is large. Gaussian isoperimetry shows that the
posterior lies near such a bulk, with the remaining signal having energy
close to a target witness's. Removing the bulk then preserves all Boolean
signs under an explicit product-norm condition. No additional restriction of prior
support is used.

Write $\sigma^2=t^2/N$, so the original Gaussian has independent entry
variances $N\sigma^2/d,\sigma^2,\sigma^2/N$. Its energy coordinates are
\[
 A_1=\sqrt{d/N}\,W^1,\qquad A_s=W^s\ (2\le s<L),
 \qquad A_L=\sqrt N\,v,
\]
and all their entries are independent $\mathcal N(0,\sigma^2)$.
Their squared Euclidean norm is exactly $\kappa_N$.

\begin{theorem}[Unrestricted learning above exponent one]
\label{thm:unrestricted-above-one}
Fix integers $d\ge1$, $L\ge2$, $C\ge1$, and a real $K\ge L$.
Use tanh hidden activations and a linear readout. Suppose that for every
$N\ge C$ there is a width-$N$, depth-$L$ target witness satisfying
\[
 f(x)z(x)\ge1\quad\hbox{for every Boolean input},\qquad \kappa_N\le K.
\]
Put
\[
 Q_* = \left(\frac{K+2}{L}\right)^{L/2}.
\]
For every fixed $\zeta>0$, every $0<\varepsilon,\delta<1$, and every
distribution $D$ on the Boolean cube, the original Gaussian learner, conditioned only
on the original global clean margin and noiseless training consistency,
has a sample bound polynomial in
\[
 d,L,C,K,Q_*,\varepsilon^{-1},\log(1/\delta)
\]
which, with probability at least $1-\delta$ over independent training
inputs drawn from $D$ and labeled by $f$, gives Gibbs error at most
$\varepsilon$, simultaneously for all
$N\ge C$ and all $0<\sigma^2\le N^{-(1+\zeta)}$.
Thus it covers all $\gamma\ge1+\zeta$ at their actual finite Gaussian scales.
The polynomial degree may depend on $\zeta$; no uniform claim as
$\zeta\downarrow0$ is made.
\end{theorem}

We prove the statement with explicit width thresholds and sample inequalities.
The hypothesis $K\ge L$ costs nothing: averaging over the cube and using
AM--GM shows that a global unit-margin witness has
$\kappa_N\ge Ld^{1/L}\ge L$.

\paragraph{A perturbation estimate using operator norms.}
Normalize hidden states by $\sqrt N$ and the augmented input by $\sqrt d$.
The input norm is at most $\sqrt2$, all coordinate activations remain
$1$-Lipschitz and zero-preserving, and the layer maps are the $A_s$ above.
If $U$ has energy at most $K+1$, and an additive weight perturbation $B$
has operator norm at most $\eta$ in every normalized layer, telescoping
one layer at a time gives
\begin{equation}\label{eq:operator-perturbation}
 \sup_x|z_{U+B}(x)-z_U(x)|
 \le\sqrt2\,L\eta
 \left[\frac{K+1+2\eta\sqrt{L(K+1)}+L\eta^2}{L-1}\right]^{(L-1)/2}.
\end{equation}
Indeed, each summand is at most $\sqrt2\eta$ times a product of $L-1$
other layer norms. The sum of squares of all perturbed layer-norm bounds
is at most the numerator, and AM--GM bounds each such product.
Set
\[
 A=\left(\frac{K+2}{L-1}\right)^{(L-1)/2},\qquad
 \eta=\min\left\{1,\frac1{4\sqrt{L(K+1)}},
                         \frac1{8\sqrt2\,LA}\right\}.
\]
The numerator in~\eqref{eq:operator-perturbation} is at most $K+2$, so the
output change is at most $1/8$. Furthermore
$A\le\sqrt e\,Q_*$, hence $\eta^{-1}$ is polynomial in $L,K,Q_*$.

\paragraph{A symmetric typical Gaussian set.}
For completely explicit constants, put
\[
 a_d^2=2(d+1)\log2+4\log8,\qquad
 b_L^2=16\log9+8\log(8L),\qquad b=\max\{a_d,b_L\}.
\]
Let $T$ consist of bulk matrices whose first and last energy-coordinate
Frobenius norms are at most $ta_d$, and whose internal operator norms are
at most $tb_L$. Then $T$ is closed, convex and centrally symmetric, and
\begin{equation}\label{eq:typical-half}
 \mu_{N,t}(T)\ge1/2.
\end{equation}
For the first block, its squared Frobenius norm is
$(t^2/N)\chi^2_{N(d+1)}$; the readout is the same expression with $d+1$
replaced by one. The exponential moment of $\chi^2$ bounds each failure
probability by $1/8$. For an internal matrix $tG/\sqrt N$, a $1/4$-net of
the unit sphere has at most $9^N$ points, and the two-net scalar Gaussian
bound gives
\[
 \Pr\{\|tG/\sqrt N\|_{\rm op}>tb_L\}
 \le2\exp(2N\log9-Nb_L^2/8)\le\frac1{4L}.
\]
Union bounding over the $L-2$ internal blocks proves~\eqref{eq:typical-half}.
Every normalized block of a member of $T$ has operator norm at most $tb$.

\paragraph{Target mass by a Gaussian shift.}
Suppose $tb\le\eta$. For a target witness $w_*$, the shifted set $w_*+T$
has target margin at least $7/8$, by~\eqref{eq:operator-perturbation}.
Central symmetry and the Gaussian change-of-mean identity yield
\begin{align*}
 \mu(w_*+T)
 &=e^{-\kappa_N(w_*)/(2\sigma^2)}
   \int_T e^{-\langle w_*,B\rangle_\kappa/\sigma^2}\,d\mu(B)\\
 &\ge e^{-K/(2\sigma^2)}\mu(T)
 \ge\tfrac12 e^{-K/(2\sigma^2)}.
\end{align*}
The inequality follows by pairing $B$ with $-B$ and using $\cosh u\ge1$.
This lower bound holds for the raw probability of the conditioning event
for every correctly labeled sample.

\paragraph{The posterior is close to the typical set.}
Let $d_\kappa(w,T)$ be Euclidean distance in energy coordinates. Gaussian
isoperimetry~\cite{borell1975} and~\eqref{eq:typical-half} imply
\[
 \mu\{d_\kappa(w,T)>r\}\le e^{-r^2/(2\sigma^2)}.
\]
Dividing by the preceding target-mass lower bound gives, for every sample,
\begin{equation}\label{eq:posterior-distance-tail}
 Q_s\{d_\kappa(w,T)^2>K+1\}\le2e^{-1/(2\sigma^2)}.
\end{equation}
On the complement, let $B$ be the metric projection of $w$ onto $T$ and
write $U=w-B$. The projection is well defined and continuous, since $T$
is nonempty, closed and convex. Then $\kappa_N(U)\le K+1$ and
$\sup_x|z_w(x)-z_U(x)|\le1/8$. The original posterior has clean margin
$>1/2$ everywhere, so $U$ has \emph{exactly the same Boolean function}
as $w$ and training margins $>3/8$. It is enough to use a margin-$1/4$
ramp loss below.

\paragraph{A polynomial width split.}
Take $\beta=\varepsilon/2$ and
\[
 M=\left\lceil\max\left\{C,2,(b/\eta)^{2/\zeta},
                                      2\log(2/\beta)\right\}\right\rceil.
\]
For every $N\ge M$ and $\sigma^2\le N^{-(1+\zeta)}$, one has
$t\le N^{-\zeta/2}$, hence $tb\le\eta$, and
$2e^{-1/(2\sigma^2)}\le\beta$.
The product norm of $U$ is at most
$d^{-1/2}[(K+1)/L]^{L/2}$. Write
\[
 G=\sqrt{(d+1)/d}\left(\frac{K+1}{L}\right)^{L/2},\qquad H=L-1.
\]
The product-norm Rademacher bound, now with the $4$-Lipschitz ramp loss
$\min\{1,\max\{0,1-4yz\}\}$, gives one sample event of probability at
least $1-\delta/2$ on which every such consistent residual network has
error at most
\begin{equation}\label{eq:distance-sample-bound}
 F_n=\frac{8G}{\sqrt n}\left[2+\sqrt{2H\log(3+nH^2)}\right]
                     +\sqrt{\frac{\log(2/\delta)}{2n}}.
\end{equation}
This event already includes all finite widths. When $F_n\le\varepsilon/2$,
adding the posterior exceptional mass $\beta$ proves the required bound
for every $N\ge M$ and every stated scale.

For $C\le N\le M$, use Theorem~\ref{thm:finitewidth-unrestricted}.
Here is a target-mass constant requiring only the present witness assumption.
Every witness weight is at most $R=\sqrt{MK}$. Put
\begin{align*}
 a&=[4(d+1+M(L-1))(MR)^{L-1}]^{-1},\\
 D_M&=\max\{1,K+[d+2+(L-2)M][2+\log(1/a)]\}.
\end{align*}
The same Gaussian-coordinate box calculation as in the finite-width proof
has physical perturbation at most $a$, retains target margin at least $3/4$,
and has mass at least $\exp(-ND_M/t^2)$.
Applying that theorem with error $\varepsilon$ and failure probability
$\delta/2$ controls all these widths and all $t\le1$ simultaneously.
Its finite-grid logarithm is polynomial in $M,d,L,\log K$.
The union of the two sample events proves the theorem.
Finally, $G\le\sqrt2 Q_*$, $b^2=O(d+\log L)$ and
$A\le\sqrt e Q_*$ show that $M$ and a sample size satisfying all the
displayed inequalities are polynomial in the stated arguments for fixed
$\zeta$. This proves the asserted quantitative dependence.

\paragraph{Scope.}
This argument does not apply at $\gamma=1$: the Gaussian bulk's operator
norm does not tend to zero with width there. It also does not make $Q_*$
polynomial for every small circuit. If a fixed base witness can be padded
with energy $K=L+B$, where $B$ is independent of the added depth, then
$Q_*\le e^{(B+2)/2}$. The theorem therefore removes exponential dependence
on \emph{additional depth}, at every width and every fixed exponent above
one, while retaining a potentially exponential dependence on base-witness
energy. Stronger target-specific energy estimates can make $Q_*$ itself
polynomial.

\section{A polynomially quantified mean-field prefactor window}
\label{sec:small-prefactor}

The bulk-subtraction proof depends on the prefactor $t$, not directly
on the exponent $\gamma$. This gives a separate finite-temperature
statement at mean-field scaling. It does not include the original
critical convention $t=1$.

\begin{theorem}[Small prefactor at every width]
\label{thm:small-meanfield-prefactor}
Assume the witness hypothesis of
Theorem~\ref{thm:unrestricted-above-one}, with parameters $d,L,C,K$
and $Q_*=[(K+2)/L]^{L/2}$. Define
\[
 A=\left(\frac{K+2}{L-1}\right)^{(L-1)/2},\qquad
 \eta=\min\left\{1,\frac1{4\sqrt{L(K+1)}},
                         \frac1{8\sqrt2\,LA}\right\},
\]
\[
 b^2=\max\{2(d+1)\log2+4\log8,\;16\log9+8\log(8L)\},
\qquad
 t_0=\min\left\{1,\frac\eta b,
                    \frac1{\sqrt{2\log(4/\varepsilon)}}\right\}.
\]
There is a sample bound polynomial in
$d,L,K,Q_*,\varepsilon^{-1},\log(2/\delta)$ such that, with probability
at least $1-\delta$, the unrestricted posterior has Gibbs error at most
$\varepsilon$, simultaneously for all $N\ge C$ and all
\[
 0<t\le t_0,\qquad \sigma^2=t^2/N.
\]
The inverse threshold $t_0^{-1}$ is polynomial in
$d,L,K,Q_*,\varepsilon^{-1}$. Neither bound depends on the realized
training sample or on its size through a circular condition.
\end{theorem}

\begin{proof}
Use the typical set $T$ and the perturbation estimate from
Section~\ref{sec:gaussian-distance}. For every width in the statement,
$tb\le\eta$, so the target shift has raw mass at least
$\frac12e^{-K/(2\sigma^2)}$. The posterior distance tail is therefore
\[
 Q_S\{d_\kappa(w,T)^2>K+1\}
 \le2e^{-1/(2\sigma^2)}
 =2e^{-N/(2t^2)}\le\varepsilon/2.
\]
Every remaining draw has exactly the same Boolean function as a
residual of energy at most $K+1$, with training margin greater than
$3/8$. The single width-independent generalization event underlying
\eqref{eq:distance-sample-bound} controls all of these residuals
simultaneously. Choosing its right-hand side at most $\varepsilon/2$
proves the claim. This sample choice is polynomial: for example, if
$G=\sqrt{(d+1)/d}[(K+1)/L]^{L/2}$, then a sufficiently large universal
constant times
\[
 [1+G+L+\varepsilon^{-1}+\log(4/\delta)]^{12}
\]
suffices. Finally $G\le\sqrt2Q_*$ and $A\le\sqrt eQ_*$ give all the
stated polynomial dependences. This proof requires no width split,
because the small-bulk condition has been imposed directly on $t$.
\end{proof}

\paragraph{A sample bound without an additional depth factor.}
The global-margin counting theorem gives a sharper choice of sample
size in Theorem~\ref{thm:small-meanfield-prefactor}. Put
\[
 \overline P=\max\left\{1,\frac8{3\sqrt d}
                  \left(\frac{K+1}{L}\right)^{L/2}\right\}.
\]
Rescale the residual readout by $8/3$. Its Boolean function now belongs
to $\mathcal H_{\overline P}$ from
Corollary~\ref{cor:all-depth-pf-counting}. Therefore it suffices to use
\[
 n\ge\frac2\varepsilon
   \left[p(\overline P)+\overline P^2\log(d+1)
                         +\log(2/\delta)\right],
\]
for a universal polynomial $p$. In particular, if $K\le L+B$ with
$B$ independent of the learner depth, then
\[
 \overline P\le\max\{1,(8/3)d^{-1/2}e^{(B+1)/2}\}.
\]
The required sample size has no remaining learner-depth dependence.
The sufficient prefactor $t_0$ still depends on $L$, and the original
critical choice $t=1$ is still excluded. Likewise, when this argument
is combined with the small-width fallback below, that finite-width
part retains its polynomial dependence on $L$.

\begin{corollary}[Polynomial product targets at mean-field scaling]
\label{cor:small-prefactor-product-targets}
Suppose $f$ has an arbitrary finite-width, depth-$D_0$ tanh
representation with unit global margin and product norm at most
$P\ge1$. There are universal polynomials $p,q$ and a universal
constant $c$ such that, whenever
\[
 L\ge D_0+c\log(2+d+D_0+P)
 \quad\hbox{or}\quad L\ge2D_0,
\]
sample size at least
\[
 p(d,D_0,P,L,\varepsilon^{-1},\log(2/\delta))
\]
gives Gibbs error at most $\varepsilon$ with probability at least
$1-\delta$, simultaneously for every actual width $N\ge C_L(f)$
and every mean-field scale
\[
 \sigma^2=t^2/N,
 \qquad 0<t\le
 q(d,D_0,P,L,\varepsilon^{-1},\log(2/\delta))^{-1}.
\]
For a bounded shallow target, or for a Boolean halfspace using
Theorem~\ref{thm:halfspace-unrestricted-all-depths}, the corresponding
statement holds at every $L\ge2$, with the same polynomial dependence
as in those positive theorems. The original critical prior $t=1$
is not asserted to satisfy this condition.
\end{corollary}

\begin{proof}
The width-compression and product-target constructions provide a
polynomial threshold $M$ and, at all $N\ge M$, unit-margin witnesses
with $K$ and $Q_*$ polynomial in $d,D_0,P,L$. Apply
Theorem~\ref{thm:small-meanfield-prefactor} to these widths. The physical
weight bound in that construction gives $C_L(f)\le M$.

For the remaining interval $C_L(f)\le N<M$, choose an actual
minimum-width bounded-weight witness and double its readout. It now
has unit margin and weights at most $2M$. Balanced cloning gives
unit-margin witnesses throughout this finite width interval with
energy polynomial in $d,L,M$. The finite-width theorem, with its
Gaussian-coordinate-box denominator, gives a polynomial sample bound
uniformly for all $t\le1$. Its constants can be bounded using $M$
and $2M$, independently of the actual minimum width. Union bound
this sample event and the large-width event, each with failure
probability $\delta/2$. Their thresholds are deterministic
polynomials of the displayed parameters. The same proof applies to
the bounded shallow and halfspace witnesses at all $L\ge2$.
\end{proof}

\section{Polynomial unrestricted learning of shallow targets at every depth}
\label{sec:shallow-unrestricted}

A shallow target can have a small first-layer Gaussian energy at a
polynomial base width simply by adding zero neurons. Its readout energy
is then polynomial. Distributing that readout cost over the learner's
remaining layers makes the product-norm parameter polynomial uniformly
in depth. This gives a family for which neither arbitrary width nor
arbitrary added depth requires a robustness restriction.

\begin{theorem}[Shallow targets, all depths and all admissible widths]
\label{thm:shallow-all-depths}
Suppose a Boolean target $f:\{-1,1\}^d\to\{-1,1\}$ has a one-hidden-layer
representation
\[
 z(x)=\sum_{j=1}^{C}v_j\tanh(a_j^T\bar x),\qquad
 \bar x=(1,x),\qquad f(x)z(x)\ge1
\]
with every coordinate of $a_j$ and every $v_j$ bounded in absolute value
by $R\ge1$. Let $C_L(f)$ denote its bounded-weight network complexity
at depth $L$: the least width $N$ admitting a realization with all
weights at most $N$ in absolute value and signed output greater than
$1/2$ everywhere.

For each fixed $\zeta>0$, the original unrestricted Gaussian learner has
sample complexity polynomial in
\[
 d,C,R,L,\varepsilon^{-1},\log(1/\delta),
\]
for every specified learner depth $L\ge2$ and every input distribution,
simultaneously for every width $N\ge C_L(f)$ and every
$\gamma\ge1+\zeta$. With probability at least
$1-\delta$ over a sample of that size, posterior Gibbs error is at most
$\varepsilon$ simultaneously at all these widths and all prior scales
$0<\sigma^2\le N^{-(1+\zeta)}$. The polynomial's degree may depend on
$\zeta$. The prior is conditioned only on the original global clean
margin and the training labels.
\end{theorem}

\begin{proof}
Let $A$ have rows $a_j^T$, and set
\[
 M=\left\lceil\max\{4,C,2R,d\|A\|_F^2\}\right\rceil,
 \qquad T=\max\{1,M\|v\|_2^2\}.
\]
Both are polynomial in $d,C,R$. Add zero hidden neurons to obtain width
$M$. This preserves the function, and its first-layer and readout
energies satisfy
\[
 E_{\rm pre}=\frac dM\|A\|_F^2\le1,
 \qquad c^2=M\|v\|_2^2\le T.
\]
In fact $M\|v\|_2^2\ge1$, since unit margin and $|\tanh|\le1$ imply
$\|v\|_2\ge1/\sqrt C$, but the maximum makes this unneeded.

For any $N\ge M$, balanced cloning of this shallow network preserves its
output exactly and gives
\[
 E_{\rm pre}\le2,\qquad c^2=N\|v_N\|_2^2\le2T.
\]
Explicitly, assign each of the $M$ original neurons either
$\lfloor N/M\rfloor$ or $\lceil N/M\rceil$ clones, copy its incoming
weights, and divide its readout weight by its number of clones.
The first cost increases by at most $M\lceil N/M\rceil/N\le2$,
and the readout cost by at most $N/(M\lfloor N/M\rfloor)\le2$.
For learner depth $L=2$ we already have a unit-margin witness with
\[
 K_2=2+2T,\qquad
 Q_*^2=\left(\frac{K_2+2}{2}\right)^2=(T+2)^2.
\]

For $L\ge3$, apply Lemma~\ref{lem:scalar-transmission} with
$D=4c^2$ and $m=L-1$. Replace the old readout by a layer whose every
row is $2\alpha v_N^T$. Follow it with the scalar transmission layers,
implemented by equal coordinates with dense matrix entries $b_s/N$,
and final readout entries $q/N$, where
\[
 q^2=\max\{c^2,1\}^{1/(L-1)}\le(2T)^{1/(L-1)}.
\]
The target had margin at least one, hence in particular at least
$1/2$, so the scalar construction gives margin at least $1/2$.
Double its final readout to restore margin at least one. This adds
$3q^2$ to the energy. Thus every width $N\ge M$ has a depth-$L$
unit-margin witness with
\[
 \kappa_N\le K_L:=2+(L+14)(2T)^{1/(L-1)}.
\]
Writing $q_*^2=(2T)^{1/(L-1)}\ge1$, we have
$K_L+2\le(L+18)q_*^2$. Consequently
\begin{equation}\label{eq:shallow-product-polynomial}
 Q_*^2=\left(\frac{K_L+2}{L}\right)^L
 \le e^{18}(2T)^{L/(L-1)}
 \le e^{18}(2T)^{3/2}\qquad(L\ge3).
\end{equation}
The exponent $L/(L-1)$ is at most $3/2$ for $L\ge3$.
Theorem~\ref{thm:unrestricted-above-one}, with base width $M$, therefore
provides a sample bound polynomial in the stated quantities for all
$N\ge M$.

We verify admissible smaller widths rather than silently identifying
$M$ with the target's minimum complexity. First, $C_L(f)\le M$ for all
$L\ge2$. At width $M$, the original incoming weights have magnitude
at most $R\le M/2$. In the added scalar channel,
$\alpha\le\operatorname{atanh}(1/2)<0.55$, so its first weights
$2\alpha v_j$ have magnitude at most $1.1R<M$.
The scalar construction has $b_s\le2.2q$ and final weights $2q/M$.
Since $T\le MC R^2\le M^4$ and $L\ge3$,
\[
 q\le(2T)^{1/4}\le2^{1/4}M,
 \qquad b_s/M<3\le M,\qquad 2q/M<3\le M.
\]
Thus this depth-$L$, width-$M$ circuit satisfies the bounded-weight
convention and has margin at least one. The case $L=2$ is immediate.

For $C_L(f)\le N<M$, start from a defining width-$C_L(f)$ circuit,
clone it to width $N$, and double its readout. It has margin greater
than one, all weights at most $2M$, and Gaussian energy bounded by a
polynomial in $d,L,M$. The Gaussian box estimate and
Theorem~\ref{thm:finitewidth-unrestricted} consequently give a sample
bound polynomial in $d,L,M,\varepsilon^{-1},\log(1/\delta)$,
simultaneously for all these widths and all $t\le1$.
Use failure probability $\delta/2$ for this event and $\delta/2$ for
the preceding event for $N\ge M$. Their intersection proves the claim.
\end{proof}

The complexity parameter controlling this theorem's sample size is the
size and weight bound of a \emph{shallow} representation. Its width
admissibility is the actual depth-$L$ complexity $C_L(f)$. The theorem
does not identify shallow complexity with the minimum complexity of an
arbitrary deep circuit, and therefore does not settle the unrestricted
problem for all small deep circuits. It also makes no claim at $\gamma=1$.

\begin{corollary}[Parity at arbitrary depth and admissible width]
\label{cor:parity-full-polynomial}
For every fixed $\zeta>0$, parity of any $k\ge1$ chosen coordinates in
$\{-1,1\}^d$ has unrestricted Gaussian posterior sample complexity
polynomial in $d,k,L,\varepsilon^{-1},\log(1/\delta)$, for every
$L\ge2$, every $N\ge C_L(f)$, and every $\gamma\ge1+\zeta$.
The bound holds for every input distribution and for the original
finite-variance posterior.
\end{corollary}

\begin{proof}
Let $q=(k-\sum_{i\in S}x_i)/2\in\{0,\ldots,k\}$ count negative
coordinates, and put $a=\log(16k)$. Define
\[
 \tau_j(x)=\tanh\bigl(a(q-j+1/2)\bigr),\qquad 1\le j\le k.
\]
Every argument has magnitude at least $a/2$, so
$|\tau_j-\operatorname{sign}(q-j+1/2)|\le2e^{-a}=1/(8k)$.
Writing $c_k=1$ for even $k$ and $c_k=0$ for odd $k$, the exact step
identity is
\[
 (-1)^q=c_k+\sum_{j=1}^k(-1)^j
                  \operatorname{sign}(q-j+1/2).
\]
Hence $F=c_k+\sum_j(-1)^j\tau_j$ satisfies
$|F-(-1)^q|\le1/8$. Represent the constant by one extra hidden neuron
with augmented-input weight $\operatorname{atanh}(1/2)$ and no other
incoming weights. A linear readout equal to $2F$ has signed margin
at least $7/4$. Its threshold-neuron readout weights are $2(-1)^j$,
and its constant-neuron weight is $4c_k$.
Its incoming weights have magnitude at most $a(k+1)$.
Thus $C=k+1$ and $R=\max\{4,a(k+1)\}$ are polynomial in $k$.
Apply Theorem~\ref{thm:shallow-all-depths}.
\end{proof}

\begin{corollary}[Symmetric Boolean targets]
The conclusion of Corollary~\ref{cor:parity-full-polynomial} holds for
any Boolean function of the Hamming weight of $k$ selected coordinates.
\end{corollary}

\begin{proof}
If its values are $h(0),\ldots,h(k)\in\{-1,1\}$, then
\[
 h(q)=\frac{h(0)+h(k)}2+
       \sum_{j=1}^k\frac{h(j)-h(j-1)}2
                    \operatorname{sign}(q-j+1/2).
\]
Each coefficient has magnitude at most one. Replacing the signs by
the same $\tau_j$ and doubling the output gives margin at least $7/4$,
using the same width and weight bounds as in the parity proof.
\end{proof}


\begin{corollary}[Polynomial-size DNF and CNF targets]
A Boolean function on $d$ bits represented by a DNF or CNF with $s\ge1$
clauses has unrestricted Gaussian posterior sample complexity polynomial
in $d,s,L,\varepsilon^{-1},\log(1/\delta)$, for every specified $L\ge2$,
every $N\ge C_L(f)$, and every $\gamma\ge1+\zeta$, with $\zeta>0$ fixed.
\end{corollary}

\begin{proof}
For a DNF, let $q_j(x)$ be the number of false literals in clause $j$,
and put $a=\log(16s)$ and
$\tau_j(x)=\tanh(a(1/2-q_j(x)))$.
The exact clause indicator is
$[1+\operatorname{sign}(1/2-q_j)]/2$.
Thus
\[
 H(x)=s-1+\sum_{j=1}^s\operatorname{sign}(1/2-q_j(x))
\]
is $-1$ when no clause is satisfied and at least $1$ otherwise.
Replacing the signs by $\tau_j$ changes $H$ by at most $1/8$.
Doubling the resulting expression gives global signed margin at least
$7/4$. It is a shallow tanh network with $s+1$ features: the extra
feature is the constant $1/2$, with readout weight $4(s-1)$.
After removing duplicate literals and discarding contradictory clauses,
each clause has at most $d$ literals. Each $q_j$ is affine in the bits,
so the remaining incoming and readout weights have magnitude at most
$R=\max\{4s,a(d+1)\}$. The shallow-target theorem applies.
A CNF is the negation of a DNF in negated literals, and negating the
readout preserves every bound. Constant functions are immediate.
\end{proof}

\subsection{Targets of any fixed intrinsic depth}
\label{sec:fixed-intrinsic-depth}

The shallow-target argument extends to any fixed target depth by using
different numbers of clones at its successive hidden layers. The degree
of the resulting polynomial may depend on the target depth, but not on
the learner's additional depth.

\begin{corollary}[Fixed target depth, arbitrary learner depth]
\label{cor:fixed-intrinsic-depth}
Fix an integer $D_0\ge2$ and $\zeta>0$. Suppose $f$ has a width-$C$,
depth-$D_0$ tanh realization in the present architecture, with signed
output at least one on the Boolean cube and all weights bounded by
$R\ge1$. For every specified learner depth $L\ge D_0$, the original
unrestricted Gaussian learner has sample complexity polynomial in
\[
 d,C,R,L,\varepsilon^{-1},\log(1/\delta).
\]
The degree and constants may depend on $D_0$ and $\zeta$.
For every input distribution, one sample event of probability at least
$1-\delta$ gives Gibbs error at most $\varepsilon$ simultaneously for
every $N\ge C_L(f)$ and every
$0<\sigma^2\le N^{-(1+\zeta)}$.
Here $C_L(f)$ is the actual bounded-weight network complexity defined in
Theorem~\ref{thm:shallow-all-depths}. The prior retains the original
global clean-margin conditioning and sign-only training likelihood.
\end{corollary}

\begin{proof}
Put $H=D_0-1$. Pad the target's hidden layers to width $C$ with zero
neurons if necessary, and write $W^1,\ldots,W^H,v$ for its weights.
Choose an integer
\[
 B=\left\lceil\max\left\{4,2R,\frac dC\|W^1\|_F^2,
                 \|W^2\|_F^2,\ldots,\|W^H\|_F^2\right\}\right\rceil,
 \qquad M=CB^H,
\]
omitting the internal-layer terms when $H=1$.
Since $\|W^1\|_F^2\le C(d+1)R^2$ and
$\|W^s\|_F^2\le C^2R^2$ internally, $B$ is polynomial in $d,C,R$,
and so is $M$ for fixed $D_0$.

At hidden layer $s$, give each original neuron exactly
\[
 n_s=B^{H-s}
\]
active copies, and fill the other positions up to width $M$ with zero
neurons. Copy first-layer rows. For $s\ge2$, give every edge from a copy
of source $j$ to a copy of target $i$ weight $W^s_{ij}/n_{s-1}$.
All unused edges are zero. Since $n_H=1$, retain the original readout
on the final active neurons. This preserves the function exactly.
The first-layer and internal energy terms are respectively
\[
 \frac dM n_1\|W^1\|_F^2
      =\frac d{CB}\|W^1\|_F^2\le1,
 \qquad
 \frac{n_s}{n_{s-1}}\|W^s\|_F^2
      =\frac1B\|W^s\|_F^2\le1.
\]
Thus the energy before the readout is at most $H$. Put
\[
 T=\max\{1,M\|v\|_2^2\}\le MC R^2.
\]
It is polynomial in $d,C,R$ for fixed $D_0$.

Balanced cloning of this width-$M$ network into any $N\ge M$ preserves
the function and gives
\[
 E_{\rm pre}\le2H,\qquad c^2=N\|v_N\|_2^2\le2T.
\]
If $L=D_0$, these are already unit-margin witnesses with
$K=2H+2T$, and
\[
 Q_*^2=\left(\frac{K+2}{D_0}\right)^{D_0}
       \le(4T)^{D_0}.
\]

If $L>D_0$, set $m=L-H\ge2$. Replace the old readout by the scalar
transmission from Lemma~\ref{lem:scalar-transmission}, using $D=4c^2$,
and double the final scalar readout to obtain unit margin. Its readout
cost before doubling is
$q^2=\max\{c^2,1\}^{1/m}\le(2T)^{1/m}$.
The total energy is at most
\[
 K_L=2H+(m+15)(2T)^{1/m}.
\]
This is the desired depth: the $H$ old hidden layers and $m-1$ new hidden
layers total $L-1$. Since $(2T)^{1/m}\ge1$,
\begin{align*}
 Q_*^2
 &=\left(\frac{K_L+2}{L}\right)^L\\
 &\le\left(1+\frac{H+17}{L}\right)^L(2T)^{L/(L-H)}\\
 &\le e^{H+17}(2T)^{(D_0+1)/2}.
\end{align*}
The last step uses $L/(L-H)=1+H/m\le1+H/2$.
For $L\ge2D_0$, the sharper upper bound
$Q_*^2\le e^{H+17}(2T)^2$ also holds.
Both $K_L$ and $Q_*$ are polynomial in the stated parameters when
$D_0$ is fixed, uniformly over the learner depth.
Theorem~\ref{thm:unrestricted-above-one} now controls every $N\ge M$.

It remains to check admissible widths below $M$.
First $C_L(f)\le M$ at every $L\ge D_0$.
The tapered construction has all original weights at most $R\le M/2$.
In a padded scalar channel, $\alpha\le\operatorname{atanh}(1/2)<0.55$,
so its first added weights $2\alpha v_j$ have magnitude less than $M$.
The scalar construction has $b_s<2.2q$, and
\[
 T\le MC R^2\le M^4,\qquad
 q\le(2T)^{1/4}\le2^{1/4}M.
\]
Its dense internal entries $b_s/M$ and final readout entries $2q/M$
are therefore less than $3\le M$.
This gives a bounded-weight width-$M$ witness with unit margin at each
learner depth; for $L=D_0$ no scalar channel is needed.

For $C_L(f)\le N<M$, clone a defining width-$C_L(f)$ realization to
width $N$ and double its readout. Its margin is greater than one and
all weights are at most $2M$. The Gaussian box estimate and
Theorem~\ref{thm:finitewidth-unrestricted} give a bound polynomial in
$d,L,M,\varepsilon^{-1},\log(1/\delta)$ for all these widths and all
$t\le1$. Allocate failure probability $\delta/2$ to this event and to
the preceding event for $N\ge M$. Their intersection proves the claim.
\end{proof}

Thus polynomial dependence on learner depth holds for targets of any
fixed intrinsic depth with polynomial target-size dependence. The width
$M=CB^{D_0-1}$ and the exponents above explain why this does not give a
single polynomial jointly in the intrinsic target depth and target size.
The mean-field exponent $\gamma=1$ is also still outside this result.

\begin{corollary}[Quantitative dependence on intrinsic depth]
\label{cor:intrinsic-depth-quantitative}
Fix $\zeta>0$. There are constants $a_\zeta,b_\zeta>0$, independent of
the intrinsic depth, with the following property. Suppose $f$ has a
width-$C$, depth-$D_0$ tanh realization with unit signed margin and
weights bounded by $R\ge1$, where $D_0\ge2$, and set
$S=2+d+C+R$. For every specified learner depth $L\ge D_0$, the sample
bound
\[
 n\ge\left\lceil S^{a_\zeta D_0}
       \bigl(L+\varepsilon^{-1}+\log(2/\delta)\bigr)^{b_\zeta}
       \right\rceil
\]
suffices for posterior Gibbs error at most $\varepsilon$, with probability
at least $1-\delta$ for every input distribution, simultaneously for all
$N\ge C_L(f)$ and $0<\sigma^2\le N^{-(1+\zeta)}$.
The prior and likelihood are those of
Corollary~\ref{cor:fixed-intrinsic-depth}.

Thus the dependence on intrinsic depth is
$\exp[O_\zeta(D_0\log S)]$, while the dependence on learner depth remains
polynomial with degree independent of $D_0$. In particular, if
$D_0\le c\log S$ for a fixed constant $c$, the target-parameter factor
is at most $\exp[O_{\zeta,c}((\log S)^2)]$.
This is a quasipolynomial bound in $S$, not a single polynomial jointly
in intrinsic depth and target size.
\end{corollary}

\begin{proof}
First suppose $L\ge2D_0$. Use the tapered construction in the preceding proof, with
$H=D_0-1$, $M=CB^H$, and $T=\max\{1,M\|v\|_2^2\}$.
Its explicit definition gives $B\le S^5$, so
$M\le S^{5D_0}$ and $T\le MC R^2\le S^{7D_0}$.
Put $m=L-H$. Since $L\ge2D_0$, one has $L/m\le2$ and
\[
 K_L=2H+(m+15)(2T)^{1/m}=O(LT),\qquad
 Q_*^2\le e^{H+17}(2T)^{L/m}
          \le e^{H+17}(2T)^2=S^{O(D_0)}.
\]

For the remaining range $D_0\le L<2D_0$, do not taper the target.
Instead pad its original realization to
$M_0=\lceil\max\{4,C,2R\}\rceil\le S^2$, and put
$T_0=\max\{1,M_0\|v\|_2^2\}\le S^5$.
Balanced cloning to every $N\ge M_0$ gives
\[
 E_{\rm pre}\le2H S^4,\qquad c^2\le2T_0.
\]
At $L=D_0$, retain the readout. For $L>D_0$, append the same scalar
transmission with $m=L-H\ge2$. These constructions give, respectively,
\[
 K_L\le2H S^4+2T_0,
 \qquad
 K_L\le2H S^4+(m+15)(2T_0)^{1/m}.
\]
In either case $K_L\le19LS^5$. Therefore
\[
 Q_*^2\le(20S^5)^L\le S^{7L}\le S^{14D_0}.
\]
The width-$M_0$ witness has weights bounded by $M_0$: the original
weights satisfy $R\le M_0/2$, and the scalar-channel check from the
preceding proof applies because $T_0\le M_0^4$.
Thus $C_L(f)\le M_0$ in this range.

In both depth ranges, the polynomial in
Theorem~\ref{thm:unrestricted-above-one} has degree and constants
depending only on $\zeta$. Substituting the corresponding bounds gives
the stated sample form above the respective base width $M$ or $M_0$.
For smaller admissible widths, the defining-witness and Gaussian-box
argument in the preceding proof gives a sample bound polynomial in
$d,L,M,\varepsilon^{-1},\log(1/\delta)$ (or with $M_0$ in place of $M$),
with degree independent of $D_0$. Substitution gives the same sample form.
Split the failure probability between the two width ranges and enlarge
$a_\zeta,b_\zeta$ to finish.
\end{proof}

\section{A finite-sample extension to targets close to a shallow function}
\label{sec:unrestricted-target-corruption}

The denominator witness need only interpolate the observed labels. This
permits a limited positive extension beyond targets having an exact cheap
shallow representation. The restriction below is quantitative: the target
disagreement probability times the chosen sample size must fit a prescribed
patch budget. It is not an unrestricted theorem for every fixed nonzero
corruption rate at arbitrarily large sample sizes or confidence levels.

\begin{lemma}[Patching finitely many Boolean inputs]
\label{lem:shallow-point-patching}
Suppose a shallow width-$C$ tanh network $z_g$ has weights bounded by
$R\ge1$ and satisfies $g(x)z_g(x)\ge1$ on the Boolean cube.
For an integer $b\ge0$, define
\[
 Z=CR,\quad a=\log(16(b+1)(Z+1)),\quad
 C_b=C+2b+1,\qquad R_b=\max\{2R,4Z,(d+1)a\}.
\]
For any set $P$ of at most $b$ cube inputs, the Boolean function obtained
by flipping $g$ exactly on $P$ has a shallow width-$C_b$, weight-bound-$R_b$
representation with unit signed margin.
\end{lemma}

\begin{proof}
For each $p\in P$, let
\[
 h_p(x)=\tanh\bigl(a(p\cdot x-d+1/2)\bigr),\qquad
 I_p(x)=\frac{1+h_p(x)}2,\qquad c_p=-2z_g(p).
\]
Then $|I_p(x)-\mathbf1\{x=p\}|\le e^{-a}$ and $|c_p|\le2Z$.
The ideal expression $z_g(x)+\sum_pc_p\mathbf1\{x=p\}$ has signed
margin at least one for the flipped function. Replacing each indicator
by $I_p$ changes it by at most $2Zb e^{-a}<1/8$.
Double the resulting expression. Its original-feature readout weights
are $2v_j$; its isolator readout weights are $c_p$; and its additional
constant output is $\sum_pc_p$.
Implement this constant with $|P|+1$ neurons of activation $1/2$, each
having readout weight $2\sum_pc_p/(|P|+1)$, of magnitude at most $4Z$.
Their bias is $\operatorname{atanh}(1/2)$ and their other incoming
weights are zero. The number and weight bounds follow, and the signed
margin is greater than $7/4$.
\end{proof}

\begin{theorem}[A finite-sample target-corruption guarantee]
\label{thm:unrestricted-target-corruption}
Fix $\zeta>0$. There is a polynomial $P_\zeta$, whose degree and constants
depend only on $\zeta$, with the following property. Let $D$ be any input
distribution on the Boolean cube, let $f$ be any Boolean target, and let
$g$ have the shallow representation in Lemma~\ref{lem:shallow-point-patching}.
Put $\rho=\Pr_D\{f(X)\ne g(X)\}$.
Fix a learner depth $L\ge2$, $0<\varepsilon,\delta<1$, and an integer
patch budget $b\ge4\log(2/\delta)$.
For any integer sample size $n$ satisfying both
\[
 n\ge P_\zeta(d,L,C_b,R_b,\varepsilon^{-1},\log(2/\delta)),
 \qquad \rho n\le b/4,
\]
the original Gaussian posterior, conditioned only on global clean output
margin and on ordinary training signs, has Gibbs error at most
$\varepsilon$ with probability at least $1-\delta$.
The sample event is simultaneous for every actual admissible width
$N\ge C_L(f)$ and every $0<\sigma^2\le N^{-(1+\zeta)}$.
No shallow representation or circuit-size bound for the exact target $f$
is assumed.
\end{theorem}

\begin{proof}
The number of observed disagreements between $f$ and $g$ is binomial with
mean $\rho n\le b/4$. A Chernoff bound gives
\[
 \Pr\{\text{more than }b\text{ disagreements}\}
 \le(e/4)^b\le\delta/2.
\]
On the complementary event, flip $g$ at the distinct observed disagreement
inputs. Lemma~\ref{lem:shallow-point-patching} supplies a globally
unit-margin, sample-consistent shallow witness with the same deterministic
parameter bounds $C_b,R_b$, independently of which disagreements occur.

Apply zero-padding, balanced cloning, and scalar transmission to this
witness. With
\[
 M=\left\lceil\max\{4,C_b,2R_b,dC_b(d+1)R_b^2\}\right\rceil,
 \qquad T=MC_bR_b^2,
\]
one obtains unit-margin interpolants at every $N\ge M$, with energy
$K_2=2+2T$ at depth two and
$K_L=2+(L+14)(2T)^{1/(L-1)}$ at depth $L\ge3$.
Their energy-product parameters are polynomial in $d,C_b,R_b$, uniformly
over $L$.

The Gaussian-distance proof uses its witness only to lower-bound the raw
probability of the posterior conditioning event. A globally margin-valid
sample interpolant suffices for that step; its labels away from the sample
need not equal $f$. Therefore the same proof yields, on a sample event of
probability at least $1-\delta/2$, the desired risk bound whenever the
disagreement budget is met. The uniform Rademacher and finite-grid sample
events depend only on the displayed deterministic bounds, not on the
locations or labels of the patched inputs. There is no union bound over
patch sets in this argument.

For completeness, this also handles every smaller actual admissible width.
The Gaussian-distance argument has a deterministic polynomial width cutoff
$W\ge M$. For $C_L(f)\le N<W$, a defining width-$C_L(f)$ realization of
the true target has weights at most $C_L(f)\le W$. Clone it to width $N$
and double its readout. Its weights are at most $2W$, and its energy and
Gaussian-box target-mass exponent are polynomial in $d,L,W$.
The finite-grid proof therefore applies uniformly to this width range.
If $C_L(f)\ge W$, that range is empty. In either case its sample bound has
the stated polynomial form and does not require an upper bound on
$C_L(f)$ as a separate parameter. Combining the two sample events proves
the theorem.
\end{proof}

For example, choose $b=\lceil4\log(2/\delta)\rceil$ and take $n$ to be
the ceiling of the displayed polynomial bound. The permitted corruption
level $b/(4n)$ is inverse polynomial in the base representation and the
statistical parameters. This is stronger than the elementary comparison
that requires observing no disagreement, whose failure probability is
controlled only when $\rho n\lesssim\delta$. Here a nonzero number of
contradictory training labels is explicitly interpolated.
The two inequalities on $n$ are essential: this argument does not imply
that the guarantee persists for every larger sample size at a fixed
positive corruption probability.

\section{A limitation of neuronwise cloning, even at polynomial path weight}
\label{sec:neuronwise-cloning-barrier}

A circuit's absolute path weight is the sum, over its input-to-output
paths, of the products of absolute edge weights. Even a circuit with
polynomial absolute path weight can have a
quasipolynomial minimum product of layer Frobenius norms under \emph{all}
neuronwise cloning choices. The statement below concerns an exact
representation-preserving method; it is not a lower bound on the minimum
energy of the represented Boolean function, which may have a different
and much cheaper realization.

\begin{proposition}[Complete-tree cloning obstruction]
Let $H\ge1$, $d=2^H$, and consider a complete binary computation tree with
$H$ hidden layers. Layer $s$ contains $n_s=2^{H-s}$ computational neurons.
Each neuron has two computational incoming edges of absolute weight $a>0$;
the first layer takes its two inputs from distinct Boolean coordinates.
The final scalar readout has absolute weight $\beta>0$.
Arbitrary additional bias-supply neurons and edges are permitted.

Replace every computational neuron by any positive integer number of
identical copies, divide its outgoing edges by the number of its copies,
and fill unused positions with zero neurons to obtain common hidden width
$N$. Write $E_1,\ldots,E_{H+1}$ for the weighted squared layer norms of the
resulting network. Then
\[
 \prod_{s=1}^{H+1}E_s
 \ge d\beta^2(2a^2)^H 2^{H(H-1)/2}.
\]
Consequently its Gaussian energy satisfies
\[
 \kappa_N\ge (H+1)
 \left[d\beta^2(2a^2)^H 2^{H(H-1)/2}\right]^{1/(H+1)}.
\]
The computational input-to-output absolute path sum is merely
$\beta(2a)^H$. For fixed $a,\beta>0$, this is polynomial in $d$, whereas
the lower bound on the product of layer energies is
$\exp(\Omega((\log d)^2))$.
\end{proposition}

\begin{proof}
For a computational node $v$, let $r_v>0$ denote its number of copies
divided by $N$. Discard all non-computational edges when lower-bounding
layer energies. The remaining contributions are
\[
 F_1=2da^2\sum_{v\text{ at level }1}r_v,
 \qquad
 F_s=a^2\sum_{v\text{ at level }s}
       \sum_{u\text{ a child of }v}\frac{r_v}{r_u}
       \quad(2\le s\le H),
 \qquad
 F_{H+1}=\frac{\beta^2}{r_{\rm root}}.
\]
Each $F_s$ is a convex function of the log multiplicities
$y_v=\log r_v$, being a sum of exponentials of affine functions.
Average $y$ over the finite automorphism group of the complete binary
tree. The group preserves every $F_s$ and acts transitively on each level,
so Jensen's inequality decreases every $F_s$ simultaneously and makes
$r_v=r_s$ constant on level $s$. This averaging is used only for a lower
bound and need not give integer clone counts. The width constraints also
remain valid, by convexity of the exponential, although that fact is not
needed for the lower bound.

At these averaged log multiplicities,
\[
 F_1=2da^2 n_1r_1,\quad
 F_s=2a^2 n_s\frac{r_s}{r_{s-1}},\quad
 F_{H+1}=\frac{\beta^2}{r_H}.
\]
Their product telescopes to
$d\beta^2(2a^2)^H\prod_{s=1}^H n_s$.
Since $E_s\ge F_s$ before averaging and every $F_s$ decreases on
averaging, this lower-bounds the original product of layer energies.
Now $\prod_s n_s=2^{H(H-1)/2}$, and arithmetic--geometric mean gives the
energy bound. There are $2^H$ computational input-to-output paths, each
with absolute weight $\beta a^H$, proving the last assertion.
\end{proof}

In particular, allowing nonuniform clone counts within a layer does not
turn a polynomial absolute path bound into the polynomial energy-product
bound needed by the unrestricted Gaussian-distance theorem. A proof of
that implication would need a change of representation beyond exact
cloning, or a different posterior argument.


\section{What remains unresolved}

The necessary condition concerns functions up to approximation, not a
particular implementation. The depth-independent product normal form
proves that ordinary Boolean description size is polynomial in intrinsic
product complexity. The converse polynomial comparison is not known,
and even such a comparison would retain the learner-depth qualification
in our present sufficiency theorem. The exact-cloning obstruction above therefore
does not supply the missing function-level lower bound. Conversely,
the shallow and halfspace results do not identify shallow representation
size with arbitrary deep circuit size. A joint polynomial theorem for
every small deep target, or an explicit counterexample to it under this
posterior, remains open.

The exponent $\gamma=1$ lies outside the Gaussian-distance argument:
its typical hidden Gaussian bulk has order-one operator norm, so the
uniform subtraction estimate used here does not become small with width.
This observation is a proof limitation, not a sample lower bound.
Removing the original all-input output-margin conditioning is also
outside the present statements.

\begin{thebibliography}{9}
\bibitem{borell1975} C. Borell.
The Brunn--Minkowski inequality in Gauss space.
\emph{Inventiones Mathematicae} 30, 207--216, 1975.
\url{https://doi.org/10.1007/BF01425510}.
\bibitem{goldmann1998} M. Goldmann and M. Karpinski.
Simulating threshold circuits by majority circuits.
\emph{SIAM Journal on Computing} 27, 230--246, 1998.
\url{https://doi.org/10.1137/S0097539794274519}.
\bibitem{hesse2002} W. Hesse, E. Allender, and D. A. Mix Barrington.
Uniform constant-depth threshold circuits for division and iterated multiplication.
\emph{Journal of Computer and System Sciences} 65(4), 695--716, 2002.
\url{https://doi.org/10.1016/S0022-0000(02)00025-9}.
\bibitem{spira1971} P. M. Spira.
On time-hardware complexity tradeoffs for Boolean functions.
\emph{Proceedings of the Fourth Hawaii International Symposium on
System Sciences}, 525--527, 1971.
\url{https://ndlsearch.ndl.go.jp/en/books/R100000136-I1571135650878062336}.
\bibitem{golowich2018} N. Golowich, A. Rakhlin, and O. Shamir.
Size-independent sample complexity of neural networks.
\emph{COLT}, 2018.
\url{https://proceedings.mlr.press/v75/golowich18a.html}.
\bibitem{nakhleh2026} J. Nakhleh and R. D. Nowak.
Deep Neural Variation Spaces: A Unifying Perspective on Depth and Complexity.
2026. \url{https://arxiv.org/abs/2607.05546}.
\end{thebibliography}
\end{document}
